\section{Supplementary results}
\label{app:results}

\subsection{Mean-field analysis of the core model: limit cycle, healthy pockets and drift}
\label{app:mf-supp}

This subsection collects the mean-field analysis of the \emph{core model without the four rules}, that is, of the
short-lived control ($\alpha=1.4$, $a_{\max}=60$) and of its variants in $\alpha$, summarised in
Section~\ref{sec:mf-scope}; none of it applies to the calibrated candidates beyond what is stated there.

\paragraph{Occupancy.}
The effective weekly loss of social state at $N=1600$ in the short-lived control is 0.75, against a Poisson prediction of
0.001. With negative-binomial occupancy of dispersion $\xi$, the switch point $N^*$ drops from 3700 ($\xi=\infty$)
to 2200 ($\xi=2$) and 1700 ($\xi=1$) (Appendix~\ref{app:occ}).

\paragraph{Two-variable mean field: a limit cycle, not Lotka--Volterra.}
With a first-moment closure, $\mathbb E[s_f^2 s_m]\approx S^3$, the slow variables $N$ and $S$ obey
\begin{equation}
 \dot N = N\bigl[b(N)S^3-\mu\bigr],\qquad
 \dot S = r-\gamma D(N/|G|),\qquad S\in[0,1].
 \label{eq:mf2}
\end{equation}
The interior equilibrium is $D(N^*/|G|)=r/\gamma$, $S^*=(\mu/b(N^*))^{1/3}\approx0.32$--$0.37$.
Three facts fix the nature of the oscillations.
(i)~If $b$ were constant, \eqref{eq:mf2} would be separable, with first integral
$H=bS^4/4-\mu S-\int^N[r-\gamma D(n/|G|)]\,dn/n$: a continuum of neutral closed orbits around a centre, as in the
Lotka--Volterra system, although the vector field is not bilinear.
(ii)~The Allee factor gives $\operatorname{tr}J=N^*b'(N^*)S^{*3}>0$: the equilibrium is a weakly unstable focus,
with eigenvalues $0.0041\pm0.103i$ per week (linear period 61 weeks, amplitude $e$-folding time about 240 weeks).
(iii)~The clamping $S\in[0,1]$ is dissipative: an orbit that reaches $S=1$ while $N<N^*$ stays there until $N=N^*$
and then follows \emph{the} orbit through $(N^*,1)$, whatever its past.
Hence every trajectory that reaches $S=1$ joins in finite time a single periodic orbit. We call this cycle
\emph{superstable} because the clamp makes its Floquet multiplier zero: nearby orbits do not approach it
geometrically but merge with it. This is a sketch, not a proof that every trajectory gets there: it excludes the
equilibrium, and it would require showing that the growing spiral reaches $S=1$ before $N$ crosses the Allee
threshold; all the trajectories we integrated did (Fig.~\ref{fig:ode}). The period of the cycle, 96--133 weeks
depending on $\xi$ (Table~\ref{tab:ode}), brackets the $113\pm10$ weeks of the agent-based model with $\alpha=0$; this
is a range obtained by varying $\xi$, not a point prediction. The first peak is a transient overshoot of the
homogeneous initial condition $S_0=1$ (ratio 1.25 in the agent-based model, 1.9 in the ODE).

\begin{figure}[ht]
\centering
\includegraphics[width=\textwidth]{fig_mf_ode.pdf}
\caption{Two-variable mean field~\eqref{eq:mf2} with negative-binomial occupancy ($\xi=2$). (a)~$N(t)$ and $S(t)$
from $N_0=80$ and $S_0=1$: one overshoot followed by
a periodic limit cycle. (b)~Phase plane: the clamped system (blue) converges to the maximal orbit; without clamping
(green) the equilibrium is an unstable focus whose amplitude grows until it hits $S=0$ or $S=1$.}
\label{fig:ode}
\end{figure}

\paragraph{Homogeneous collapse versus healthy pockets.}
The two-variable model lacks demographic inertia. An age-structured mean field and a well-mixed variant of the
agent-based model (Poisson occupancy) both predict a single hump followed by deterministic extinction
(Fig.~\ref{fig:mfabm}): at the peak the whole young cohort reaches $s\approx0$ at once, recovering to viability
$s_v=R_0^{-1/3}\approx0.31$ takes $s_v/r\approx30$ weeks, and by then the cohort has left its fertile window. The
well-mixed agent-based model gives $N_{\max}=8860\pm80$, $t_{\rm ext}=160\pm12$ weeks and 40/40 extinctions; the
age-structured mean field gives $t_{\rm ext}=203$ weeks.
The oscillations of the spatial model are produced instead by heterogeneity. A minority of individuals living in
uncrowded cells, which we call \emph{healthy pockets} (10--20\% with $s>0.2$ at the peak), keeps $s\approx1$ and
repopulates the trough. Their effective fecundity exceeds $S^3$ by a factor 10--20 at the peak, so the first-moment
closure fails exactly where it matters. In the age-structured mean field a fixed fraction $\rho$ of pocket
individuals turns extinction into sustained oscillations above a sharp threshold $\rho_c\approx2.2\%$ (Poisson
occupancy) to $2.9\%$ ($\xi=2$), at which the trough equals the Allee threshold, $N_{\min}\approx N_A$. That
threshold assumes that pocket individuals experience no crowding at all ($f_P=0$), whereas the pockets of the
agent-based model are only less crowded, so the comparison is qualitative. The three regimes found when scanning
the aggregation strength $\alpha$ in the core model without rules are: (I)~homogeneous overshoot for small $\alpha$ with long-lived
mixing; (II)~a source--sink coexistence sustained by healthy pockets; (III)~collapse by aggregation, when clusters
exceed $m_c$ at low $N$ and no healthy pocket survives. Healthy pockets are an emergent property of the core model and
are unrelated to the refuges of rule R7.

\begin{figure}[ht]
\centering
\includegraphics[width=\textwidth]{fig_mf_vs_abm.pdf}
\caption{Population and mean social state for the spatial agent-based model with random motion, the well-mixed
agent-based model (Poisson occupancy), the homogeneous age-structured mean field and the same mean field with a 5\%
fraction of healthy-pocket individuals. Homogeneous descriptions collapse once; healthy pockets produce the
oscillatory regime.}
\label{fig:mfabm}
\end{figure}

\paragraph{Extinction and the absence of secular drift.}
Because the attractor is a cycle, extinction occurs at the troughs, when $N_{\min}$ falls below the Allee threshold
of an aged population. The short-lived control is bimodal. In the exploratory ensemble ($\alpha=1.4$, $n=100$), 49\% [39, 59] of
the colonies die in the first collapse ($t\le184$ weeks), and the survivors have a hazard of $7\times10^{-5}$ per
week [$2.6$, $15$]$\times10^{-5}$ (Fig.~\ref{fig:kmdrift}a); a piecewise-constant hazard is strongly preferred to a
constant one (likelihood ratio 148, $p<10^{-16}$), while the tail alone is compatible with an exponential (Weibull
shape 0.67, $p=0.28$). The independent production ensemble ($n=100$, $T=1000$) gives an early fraction
of 0.44 [0.35, 0.54] with the split at 1.5 times the median event time, 0.42--0.45 for splits between 1.25 and 2.5
times the median, and a tail hazard of $(0.5$--$1.1)\times10^{-4}$ per week; 53 of 100 colonies were still alive at
$T$. Extinction is a mixture of a deterministic-like early mode and a rare late mode, not a heavy-tailed law.
The inheritance rule $s_{\rm pup}=s_{\rm mother}$ was conjectured to produce a secular downward drift of successive
peaks. The mean field says otherwise: mothers are weighted by $s_f^2s_m$ in the birth flux, so newborns carry
$\bar s_{\rm birth}=\mathbb E[s^3]/\mathbb E[s^2]\ge S$; inheritance acts as \emph{selection}, and the clamp at $S=1$
erases memory in every deep trough. The slope of $\log N_{\rm peak}$ per cycle (first peak excluded) is
$+0.027\pm0.019$ for $\alpha=0$ (40 runs, 2000 weeks) and $+0.005$ [$-0.001$, $+0.012$] for $\alpha=1.4$ (50 runs
with at least four peaks; within-run regression $+0.001\pm0.002$); no run shows a significant negative
Mann--Kendall trend (Fig.~\ref{fig:kmdrift}c), and the preregistered production ensemble gave $+0.007$ [$-0.014$,
$+0.027$] (Wilcoxon $p=0.41$). Collapse in the short-lived control is thus \emph{contingent}: it requires the recovery number
$\Pi_R=r\,\tau_{\rm disp}/s_v$ ($\tau_{\rm disp}$ the fertile time left after crowding) to be below one and the
fraction of healthy pockets to be below $\rho_c$.

\begin{figure}[ht]
\centering
\includegraphics[width=\textwidth]{fig_km_drift.pdf}
\caption{(a)~Kaplan--Meier survival of the colony for the short-lived control ($\alpha=1.4$, 100 replicates, right-censored at
2000 weeks) with Greenwood band. (b)~Successive peak populations for ten runs with $\alpha=0$. (c)~Per-run slope of
$\log N_{\rm peak}$ versus cycle index: no secular drift.}
\label{fig:kmdrift}
\end{figure}

\paragraph{Demography and local occupancy.}
Let $\nu=N/|G|$ be the mean number of mice per cell ($|G|=900$). With the parameters of the short-lived control the logistic hazard
gives a life expectancy of 44.8 weeks, with death almost deterministically ``of old age''. A female that always
finds a mate and keeps $s=1$ produces $R_0=\tfrac12\,p_0\bar L q_0\,\tau_{\rm fert}\approx35$ daughters over her
reproductive window (70 with the long-lived demography of the candidates), so reproduction is limited by
\emph{mate encounter}, not by fecundity: a reproductive female conceives only if a reproductive male shares her
cell. Under random mixing the per-capita birth rate
\begin{equation}
 b(N)=\phi_F\,p_0\bar L q_0\,\bigl(1-e^{-\phi_M N/|G|}\bigr)
 \label{eq:allee}
\end{equation}
increases with $N$: the model has a built-in \emph{Allee effect}, with a deterministic extinction threshold
$N_A\approx61$ for $s\equiv1$. Here $\phi_M\approx\phi_F\approx0.42$ are the fractions of reproductive males and
females in the stationary age distribution; during growth the age structure is younger, so
Eq.~\eqref{eq:allee} is an approximation there. The latency of small founding groups and the failure of very small
ones follow from~\eqref{eq:allee}.
Social deterioration is driven by the local excess occupancy $(m_i-m_c)_+$. With independent positions the expected
weekly loss is $\gamma D(\nu)$, $D(\nu)=\mathbb E[(1+X-m_c)_+]$ with $X\sim\mathrm{Poisson}(\nu)$. Being a Poisson
tail expectation, $D$ is extremely convex ($D=0.0014$ at $\nu=2$, $0.10$ at $\nu^*=4.13$, $1.04$ at $\nu=7$), so
the social state behaves as a \emph{switch}: in the core model ($r/\gamma=0.1$) recovery dominates up to
$N^*=|G|\nu^*\approx3700$, and deterioration of order one per week sets in just above it. Pups are born in their
mother's cell, so the occupancy is clustered (index of dispersion 3--6 during growth for $\alpha=0$, up to 29 at the
peak for $\alpha=1.4$). A negative-binomial occupancy with dispersion $\xi$ captures this and lowers $N^*$ to 2200
for $\xi=2$; affinity-driven motion lowers $\xi$ further, which is why the peak decreases monotonically with
$\alpha$.

\paragraph{Irreversibility: rules R1 and R2.}
In the core model, recovery is an individual process that only needs time: an adult with $s_0\approx0$ returns to
$s_v$ in $(s_v-s_0)/r$ weeks, within its fertile window. Under R1 an adult is a ratchet ($\dot s\le0$), and recovery
becomes a \emph{lineage} process: only pups gain, and producing a pup requires reproduction, whose fecundity scales
as $s_f^2s_m$. Table~\ref{tab:lineage} gives the expected number of descendants that a lineage accumulates before
it reaches viability. It is the expected value of a branching process without learning from neighbours outside the
lineage, so it is a heuristic bound, not a prediction. With the long-lived demography of the candidates
($R_0\approx70$, $s_v=R_0^{-1/3}\approx0.24$) the expected number is at most $10^{-2}$ for $s_0\le0.05$, whatever the
gain per generation and the partner. It is 0.01--0.07 for $s_0=0.1$ when the partner is equally deteriorated, and up
to 0.7 when the partner is healthy, because fecundity then scales as $s^2$ in the first generation. A lineage of
$s_0\approx0.2$ can recover. The bottleneck at small $s_0$ is the fecundity of deteriorated parents, not the gain
per generation. In addition, every crowding hit on a healthy adult is permanent, so healthy pockets erode instead of
regenerating. The transplants of Section~\ref{sec:res-Q} confirm the bottleneck (animals with $s\approx0$ never
found a colony; animals with $s\approx0.36$ found one in 8\% of the trials under R1 against 95\% without it) but not
the second part of the argument: the offspring of intermediate founders settle at their parents' level instead of
collapsing.
Under R2 recovery is $r[(1-\lambda)+\lambda\bar s_V]$, with $\bar s_V$ the mean state of the neighbours. In the
homogeneous limit $S=0$ is absorbing for $\lambda=1$, and for $\lambda<1$
\begin{equation}
 T_{\rm rec}(\lambda)=\frac{1}{r\lambda}\,
 \ln\frac{(1-\lambda)+\lambda s_v}{(1-\lambda)+\lambda S_0},
 \label{eq:trec}
\end{equation}
so the colony is irreversible when $T_{\rm rec}$ exceeds the fertile time left after crowding. Unlike R1, R2 depends
on $r$, and in the spatial model a healthy cell is a focus from which $s$ spreads, so \eqref{eq:trec} gives lower
bounds for the threshold. A kinetic mean field with persistent cell disorder (Appendix~\ref{app:kin},
Table~\ref{tab:thr}, Fig.~\ref{fig:thr}) quantifies both rules for $r\in\{0.01,0.03\}$, the only rates it was run
with, and 10 seeds per cell (10 extinctions out of 10 have a Wilson interval of $[0.72,1]$): with the long-lived
demography, R1 gave extinction in 10/10 seeds for every $a_w\le24$ weeks at both rates, so within this range its
effect does not depend on $r$, and the agent-based plane $(r,a_w)$ later extended this to $r\in[0.02,0.3]$; R2 alone
extinguishes the long-lived colony when $\lambda$ exceeds a threshold that grows with $r$; the threshold
$\lambda_c$ is defined throughout as the first value of the grid $\{0,0.3,0.5,0.7,0.9,1\}$ at which at least half of
the 10 seeds go extinct, so it is $\lambda_c\le0.3$ at $r=0.01$ (8/10 seeds at $\lambda=0.3$; 10/10 from 0.5 on),
$\lambda_c\in(0.7,0.9]$ at $r=0.03$, and $\lambda_c\in(0.5,0.7]$ at $r=0.03$ with faster mixing; with the short
life span R2 needs $\lambda\ge0.9$ and fast mixing, and even then extinction is not certain (at most 8/10 seeds).

\paragraph{Thresholds of R1 and R2 in the kinetic mean field.}
In the kinetic mean field (Appendix~\ref{app:kin}), calibrated to the quasi-stationary regime of the short-lived control, R1
alone with the short life span and persistent healthy pockets reduces the rebound but does not extinguish the
colony ($P_{\rm ext}\le0.3$); faster mixing makes it lethal for $a_w\le12$, with weak dependence on $r$. In the
long-lived scenario ($a_{\max}=120$, fertility to 80 weeks) R1 gives $P_{\rm ext}=1$ and zero rebound for
$a_w\le24$ at both $r=0.01$ and $0.03$ (10/10 seeds per cell, Wilson interval $[0.72,1]$), with the critical window
$a_w^c\in(24,48)$ at $r=0.03$ and above 48 weeks at $r=0.01$ (same definition: first grid value with at least half
of the seeds extinct). R2 requires $\lambda\gtrsim0.9$ with the short life span, and in the long-lived scenario
$\lambda_c\le0.3$ ($r=0.01$; $P_{\rm ext}=0.8$ already at $\lambda=0.3$) or $\lambda_c\in(0.7,0.9]$ ($r=0.03$), in
agreement with Eq.~\eqref{eq:trec}; with faster mixing $\lambda_c\in(0.5,0.7]$ at $r=0.03$ (Fig.~\ref{fig:thr}). For
$\lambda=1$, Eq.~\eqref{eq:trec} gives $T_{\rm rec}=\ln(s_v/S_0)/r=274$--413 weeks for the measured
$S_{\min}\approx0.01$--0.02 and $r=0.01$.

\begin{figure}[ht]
\centering
\includegraphics[width=\textwidth]{fig_thresholds.pdf}
\caption{Extinction probability in the kinetic mean field (10 seeds, 1500 weeks) as a function of (a)~the critical
window $a_w$ of R1 and (b)~the social-learning weight $\lambda$ of R2, for the short ($a_{\max}=60$) and long ($a_{\max}=120$) life spans
and two recovery rates. Circles: persistent healthy pockets ($\xi_c=0.5$, $\tau_F=8$); faint squares: faster mixing
($\xi_c=1$, $\tau_F=5$).}
\label{fig:thr}
\end{figure}

\paragraph{Status of the exploratory mean-field statements.}
\begin{itemize}
  \item The kinetic mean field predicted that R1 with the long-lived demography gives $P_{\rm ext}=1$ and no rebound
    for every window $a_w\le24$ at the two rates it was run with ($r=0.01$ and $0.03$). The agent-based plane
    $(r,a_w)$ extends this to $r\in[0.02,0.3]$: $P_{\rm ext}=1$ at 48 of 48 points. Without the window, extinction
    occurs only for $r\le0.03$, and recurrent growth (A2) dominates for $r\ge0.07$ (Appendix~\ref{sec:res-planes}).
  \item That R1 blocks the recovery of transplanted adults is confirmed. That R2 with $\lambda=1$ alone blocks it is
    not: without R1, phase-D transplants recover in 0/60 trials at week 86 but in 10/60 at week 126, and animals of
    intermediate state found colonies in 50/60 trials (Appendix~\ref{sec:res-transplant}), as Eq.~\eqref{eq:trec}
    allows for $r=0.1$ ($T_{\rm rec}\approx27$ weeks).
  \item The ageing decline, $(T_{\rm ext}-t_{\rm pk})/t_{\rm pk}\approx2$ with a small spread, is confirmed (2.07
    with $t_{\rm pk}=\arg\max N$), but the ratio depends on how the peak time is defined on the plateau
    (Appendix~\ref{app:res-production}).
  \item Not tested in the production sweeps: the thresholds in $\lambda$ ($\lambda=1$ in all candidates), the effect
    of long-range mixing (the jump rule was discarded) and critical slowing down (the candidates do not oscillate).
\end{itemize}

\subsection{Production ensembles: details}
\label{app:res-production}

Table~\ref{tab:res-production-si} completes Table~\ref{tab:res-production} with the desirable patterns, the timings
of the boom and further magnitudes.

\begin{table}[ht]
\centering
\caption{Production ensembles in full (same ensembles and precision as Table~\ref{tab:res-production}): essential
patterns, transplants (founding transplants over trials, 60 replicates per arm; O12b pools both sexes; phase D at
$1.75\,t'_{\rm pk}$ of each colony, and the third row at week 95 in every arm), anti-patterns, magnitudes, desirable
patterns (fraction of passing runs), timings (medians over runs) and further magnitudes. ``PASS or MARGINAL'' for
O13b counts the runs in which the evaluator finds the isolated class later-born at least marginally.}
\label{tab:res-production-si}
\scriptsize
\setlength{\tabcolsep}{3pt}
\begin{tabular}{@{}llccccccc@{}}
\toprule
ID & Pattern & \textsf{C1} & \textsf{C1}$'$ & \textsf{C0} & $\alpha{=}0$ & \textsf{C2} & short-lived & long-lived \\
\midrule
O4 & peak below capacity, free space & 1.00 & 1.00 & 1.00 & 0.00 & 0.90 & 0.87 & 0.08 \\
O5 & natality ends at high $N$ & 0.99 & 0.85 & 1.00 & 1.00 & 0.81 & 0.25 & 0.90 \\
O6 & no rebound & 1.00 & 1.00 & 1.00 & 1.00 & 0.98 & 0.45 & 0.97 \\
O7 & extinction & 1.00 & 0.99 & 1.00 & 1.00 & 0.98 & 0.47 & 0.98 \\
O10 & adult deterioration before peak & 1.00 & 1.00 & 1.00 & 1.00 & 1.00 & 1.00 & 1.00 \\
O11 & late cohorts fail & 1.00 & 0.97 & 1.00 & 1.00 & 0.99 & 0.00 & 0.00 \\
O13 & two deteriorated classes & 1.00 & 1.00 & 0.92 & 1.00 & 1.00 & 0.00 & 0.00 \\
O13p & \quad persistent isolated class & 1.00 & 1.00 & 0.00 & 0.00 & 0.94 & 0.00 & 0.00 \\
O17 & Calhoun-like run (O17) & 0.99 & 0.83 & 0.92 & 0.00 & 0.72 & 0.00 & 0.00 \\
\midrule
O17p & \textbf{primary outcome} & 0.99 & 0.83 & 0.00 & 0.00 & 0.70 & 0.00 & 0.00 \\
 & \quad 95\% Wilson interval & 0.97--1.00 & 0.78--0.88 & 0.00--0.02 & 0.00--0.02 & 0.64--0.76 & 0.00--0.04 & 0.00--0.04 \\
\midrule
O12 & transplants found: phase B & 52/60 & 46/60 & 47/60 & -- & 49/60 & -- & -- \\
 & \quad phase D ($1.75\,t'_{\rm pk}$) & 0/60 & 0/60 & 0/60 & -- & 0/60 & -- & -- \\
 & \quad week 95, all arms & 0/60 & 0/60 & 0/60 & -- & 0/60 & -- & -- \\
O12b$^\dagger$ & \quad D + healthy partners & 0/120 & 0/120 & 0/120 & -- & 2/120 & -- & -- \\
A1--A4 & anti-patterns triggered & none & none & none & none & A2: 0.02 & A2: 0.55 & A2: 0.03 \\
\midrule
& $\rho_{\rm pk}=N_{\rm pk}/K_\gamma$ & 0.36 & 0.33 & 0.42 & 0.71 & 0.42 & 0.77 & 0.85 \\
& empty cells at peak, $f_0$ & 0.16 & 0.52 & 0.19 & 0.03 & 0.11 & 0.30 & 0.59 \\
& persistent isolated, $\Phi^{\rm pers}_{\rm BO}$ & 0.17 & 0.16 & 0.01 & 0.09 & 0.13 & 0.00 & 0.00 \\
& median $T_{\rm ext}$ (weeks, KM) & 169 & 182 & 166 & 155 & 218 & -- & 186 \\
& spread of $T_{\rm ext}$ (KM IQR/median) & 0.09 & 0.12 & 0.08 & 0.05 & 0.14 & -- & 0.16 \\
\midrule

O8 & senescent decline & 0.62 & 0.21 & 0.82 & 0.97 & 0.00 & 0.09 & 0.37 \\
O8b & slow initial decline & 1.00 & 1.00 & 1.00 & 1.00 & 1.00 & 1.00 & 1.00 \\
O9 & old, female-biased remnant & 0.45 & 0.29 & 0.44 & 0.49 & 0.30 & 0.21 & 0.37 \\
O10b & functional growth phase & 0.02 & 0.01 & 0.18 & 1.00 & 0.00 & 1.00 & 0.19 \\
O11b & few late-born females conceive & 1.00 & 1.00 & 1.00 & 1.00 & 1.00 & 1.00 & 0.97 \\
O13b & isolated class later-born & 0.00 & 0.00 & 0.00 & 0.00 & 0.00 & 0.00 & 0.00 \\
O14 & aggregation with free space & 1.00 & 1.00 & 1.00 & 0.00 & 1.00 & 1.00 & 1.00 \\
O15 & infant mortality $\to$ 1 & 0.01 & 0.00 & 0.00 & 0.01 & 0.01 & 0.00 & 0.01 \\
O15b & conceptions without viable young & 0.57 & 0.70 & 0.46 & 0.28 & 0.45 & 0.15 & 0.09 \\
O13b & \quad PASS or MARGINAL & 0.01 & 0.04 & 0.04 & 0.00 & 1.00 & 0.16 & 0.63 \\
\midrule
& $t'_{\rm pk}$ (weeks, median) & 34 & 35 & 33 & 31 & 62 & 31 & 45 \\
& $t_{\rm pk}=\arg\max N$ (median) & 55 & 61 & 53 & 44 & 80 & 35 & 66 \\
& plateau length (median) & 51 & 50 & 54 & 59 & 28 & 9 & 45 \\
& $t_{B99}$ (median) & 40 & 44 & 38 & 33 & 86 & 980 & 55 \\
& \quad 95\% CI of the KM median & 167--171 & 180--185 & 165--167 & 154--156 & 217--222 & -- & 181--190 \\
& $P_{\rm ext}(T)$ & 1.00 & 0.99 & 1.00 & 1.00 & 0.98 & 0.47 & 0.98 \\
& CV of $T_{\rm ext}$ (extinct runs) & 0.07 & 0.16 & 0.07 & 0.06 & 0.12 & 0.45 & 0.20 \\
& $(T_{\rm ext}-t_{\rm pk})/t_{\rm pk}$ (extinct) & 2.13 & 2.15 & 2.16 & 2.60 & 1.88 & 2.78 & 2.02 \\
\bottomrule

\end{tabular}
\end{table}

\paragraph{The plateau and the definition of the peak time.}
In the 200 runs of \textsf{C1} (medians, with interquartile ranges in brackets), the population first comes within
2\% of its maximum at week 34 [32, 36] and leaves that band at week 86 [85, 86], so the plateau lasts 51 [49, 53]
weeks. The maximum itself, $t_{\rm pk}=\arg\max N$, is reached at week 55 [50, 60], 21 [16, 25] weeks after the
plateau begins, and coincides with the last surviving birth in 69\% of the runs. The 99th percentile of the births
falls at $1.17\,t'_{\rm pk}$ [1.12, 1.24], and the decline ratio is $(T_{\rm ext}-t'_{\rm pk})/t'_{\rm pk}=3.97$
[3.72, 4.29] against 2.07 [1.85, 2.33] with $t_{\rm pk}$. The primary evaluator uses $t'_{\rm pk}$ and
$t_{B99}$; Appendix~\ref{app:res-rescoring} gives every pattern with both definitions.

\paragraph{The escaped run.}
One \textsf{C1}$'$ run in 200 was still alive at $T=600$: it fails O7 and is censored in the survival analysis. No
\textsf{C1} run escaped in the reference ensemble; in the ensemble at $T=700$ all 200 died.

\subsection{Survival of colonies}
\label{sec:res-survival}

Fig.~\ref{fig:res-survival} shows Kaplan--Meier estimates of colony survival; runs still alive at the horizon enter
as right-censored observations.
\begin{itemize}
  \item \textsf{C1} and \textsf{C0} die out within a narrow window. Their median extinction times are 169
    [167, 171] and 166 [165, 167] weeks, and the Weibull shapes of the collapse are 14.7 and 11.9.
  \item \textsf{C1}$'$ dies at 182 [180, 185] weeks, with a longer tail (Weibull shape 3.3 without a cure fraction,
    4.4 with $\hat\pi=0.005$, the single censored colony).
  \item \textsf{C2} dies later, at 218 [217, 222] weeks, and with more spread (KM interquartile range over median
    0.14). Four of its 200 colonies persist to $T=700$ and recovered natality (A2).
  \item The short-lived control is a mixture: 47 of 100 colonies died by $T=1000$, 44\% [35, 54] in the first
    collapse, the others in a quasi-stationary oscillation with a tail hazard of $0.7\times10^{-4}$ per week
    (Weibull shape 0.60 over all times; cure fraction $\hat\pi=0.53$).
  \item In the long-lived control 98 of 100 colonies die, at 186 [181, 190] weeks; with the
    initial condition of \textsf{C1} (null-model ensemble, $n=200$), 199 of 200 die, with a KM spread of 0.160.
\end{itemize}
The survival curves of \textsf{C1} and the short-lived control differ (log-rank $\chi^2=53$, $p=3\times10^{-13}$). A likelihood-ratio
test of the cure fraction is degenerate for \textsf{C1} (no censored colony) and is not reported; because no
\textsf{C1} colony dies before about 140 weeks, a Weibull without a location parameter mixes location and shape, so
the shape 14.7 should not be read as a hazard exponent. The Kaplan--Meier quantiles are the safer summary.
Among the single ablations, removing R1 is the only one that changes survival qualitatively (27 of 120 colonies go
extinct by week 600); removing AV delays extinction by about 18 weeks (median 185 against 167); removing R2 or R7
leaves the survival curve unchanged (170 and 167).

\begin{figure}[ht]
  \centering
  \includegraphics[width=\textwidth]{fig_res_survival.pdf}
  \caption{Survival of colonies.
  \textbf{(a)}~Kaplan--Meier curves of the candidates and of the short-lived and long-lived controls (reference ensemble for
  \textsf{C1}, \textsf{C1}$'$ and \textsf{C0}). Shaded bands are log--log Greenwood intervals for \textsf{C1} and the
  short-lived control. Censoring is at $T=600$ (\textsf{C1}, \textsf{C1}$'$, \textsf{C0}), $T=700$ (\textsf{C2}) and
  $T=1000$ (short-lived and long-lived controls).
  \textbf{(b)}~Single-rule ablations of \textsf{C1} ($n=120$, $T=600$).
  \textbf{(c)}~Phase-D law (Q4): time from the last birth to extinction against $a_{\max}$ ($n=300$, jittered). The
  solid line is the OLS fit. The dashed line is the expected maximum of $n_{\rm last}$ logistic lifetimes, with
  $n_{\rm last}=7$, the median number of births in the last ten weeks.}
  \label{fig:res-survival}
\end{figure}

\subsection{Ablations: factorial model and the minimal candidate}
\label{app:res-ablations}

A Firth logistic model of the full $2^4$ factorial (exploratory, Benjamini--Hochberg at 10\%) gives, on the logit
scale: positive main effects for R7 and R2 ($4.9\pm1.2$ each), none for R1 ($-0.5\pm1.1$, $p=0.63$) or AV
($-0.5\pm1.1$, $p=0.68$), positive interactions R1:R7 and R1:R2 ($6.6\pm1.9$) and weaker ones R2:AV and AV:R7
($3.9\pm1.8$, $p=0.036$). R1 acts only together with R2 or R7, because without them O17p is zero whether R1 is present
or not. Unlike in \textsf{C0} (Section~\ref{sec:calib-ablations}), stronger attraction does not rescue the ablation
of R1: with $\alpha=6$, O17p is 0.04.
Table~\ref{tab:res-c1p} compares \textsf{C1}$'$ with \textsf{C1} in the reference ensemble and in the factorial, and
gives its single ablations (the double ablations of \textsf{C1} that remove AV, $n=40$).

\begin{figure}[ht]
  \centering
  \includegraphics[width=\textwidth]{fig_res_ablations.pdf}
  \caption{Ablations of \textsf{C1}.
  \textbf{(a)}~Primary outcome O17p with the primary evaluator (filled, 95\% Wilson interval) and with the preregistered
  evaluator (open diamonds). Rows are \textsf{C1} and its single ablations ($n=120$), double to quadruple ablations
  ($n=40$), and implementation variants. The dotted line marks the 0.8 threshold for an essential pattern.
  \textbf{(b)}~Fraction of passing runs for each pattern (primary evaluator). NA marks a pattern that is not
  applicable in any run; with $\gamma=0$ the colonies never go extinct.}
  \label{fig:res-ablations}
\end{figure}

\begin{table}[ht]
  \centering
  \caption{The minimal candidate \textsf{C1}$'$ (= \textsf{C1} without AV). Top: reference ensemble ($n=200$,
  $T=600$); $p$ is the two-sided Fisher test against \textsf{C1}. Middle: the AV arm of the preregistered factorial
  ($n=120$); $p$ is the one-sided test of Q2 before Holm adjustment. Bottom: single ablations of \textsf{C1}$'$
  ($n=40$); $p$ is the one-sided Fisher test against the \textsf{C1}$-$AV arm of the factorial. O17p with 95\% Wilson
  interval; other columns are fractions of passing runs, and $f_0$ is the mean fraction of empty cells at the peak.}
  \label{tab:res-c1p}
  \small
  \begin{tabular}{@{}lrrlllllll@{}}
    \toprule
    Arm & $n$ & $T$ & O17p & $p$ & O5 & O11 & O13p & O8 & $f_0$ \\
    \midrule
    C1 & 200 & 600 & 0.99 [0.97, 1.00] &  & 0.99 & 1.00 & 1.00 & 0.62 & 0.16 \\
C1$'$ $=$ C1 $-$ AV & 200 & 600 & 0.83 [0.78, 0.88] & $1.0\times10^{-9}$ & 0.85 & 0.97 & 1.00 & 0.21 & 0.52 \\
\midrule
C1 (factorial) & 120 & 600 & 0.99 [0.95, 1.00] &  & 0.99 & 1.00 & 1.00 & 0.72 & 0.16 \\
C1 $-$ AV (factorial) & 120 & 600 & 0.79 [0.71, 0.85] & $1.1\times10^{-7}$ & 0.82 & 0.97 & 1.00 & 0.16 & 0.52 \\
\midrule
C1$'$ $-$ R1 & 40 & 600 & 0.00 [0.00, 0.09] & $<10^{-16}$ & 0.03 & 0.03 & 0.10 & 0.00 & 0.50 \\
C1$'$ $-$ R2 & 40 & 600 & 0.00 [0.00, 0.09] & $<10^{-16}$ & 0.90 & 0.00 & 0.97 & 0.53 & 0.51 \\
C1$'$ $-$ R7 & 40 & 600 & 0.00 [0.00, 0.09] & $<10^{-16}$ & 0.95 & 0.82 & 0.00 & 0.60 & 0.76 \\
\bottomrule

  \end{tabular}
\end{table}

\subsection{Confirmatory tests: further details and the preregistered evaluator}
\label{app:res-Q}

Table~\ref{tab:res-Qcmp} gives the confirmatory family with the primary evaluator and with the preregistered evaluator;
the latter replicates the preregistered analysis on the production runs.
\begin{itemize}
  \item \emph{Q1.} The transplanted \textsf{C1} animals are 37--64 weeks old at $1.75\,t'_{\rm pk}$ (week 60), 59--80
    at $1.75\,t_{\rm pk}$ (week 97) and 65--80 at week 95; the fertile window closes at 80 weeks. In arm $X$ the
    founding rates are 51/60, 12/60 and 5/60 (McNemar $p=4\times10^{-16}$, $2\times10^{-4}$ and $0.03$).
  \item \emph{Q4.} The law explains why O8 depends on longevity: O8 passes in 0.15, 0.30, 0.67, 0.83 and 0.92 of the
    runs for $a_{\max}=90$, 105, 120, 135 and 150, while every essential pattern is unchanged (O17p 0.93--1.00).
  \item \emph{Q5.} O13 fails only for the strongest aversion with the strongest attraction ($\kappa=0.3$, $\alpha=8$;
    $n^*\approx3.2$), where the crowded fraction falls to the 10\% threshold. The one-sided test of $b_2<0$ gives
    $p=1$ (Pearson $\phi=1.0$).
  \item \emph{Q6.} Profile-likelihood intervals: $[-0.15,-0.10]$ binomial and $[-0.20,-0.05]$ quasi-binomial
    ($\phi=3.1$) with the primary evaluator; $[0.20,0.40]$ and $[0.15,0.50]$ ($\phi=1.8$) with the preregistered
    evaluator. The grid of the profile is $\theta\in[-1,3]$ in steps of 0.05. The quasi-binomial $p$ is the one that
    enters the Holm sequence. The additive model in standardised $\log r$ and $\log a_w$ with quadratic terms beats
    the single-index model by $\Delta{\rm QAIC}=27.8$ over the 48 points and by 27.7 without the point that fails
    through the persistent class (dispersion of the additive model 1.9). With the preregistered evaluator the comparison
    goes the other way: the single-index model is preferred by $\Delta{\rm QAIC}=11$, so the change of criterion
    changes the adequate model as well as the sign of $\hat\theta$. With the primary evaluator the O17p boundary of
    the plane is nearly vertical in $r$ (at least 0.67 for $r\le0.10$ and at most 0.60 for $r\ge0.2$ at every
    window), which is why the prediction is reported as ``not supported'' rather than as a reversal.
  \item \emph{Q7.} The logistic fit of the preregistered plane places the extinction boundary at $\gamma=0$ at
    $w+0.24\,\Lambda=0.215$, that is $\hat\sigma=4.65$ and $\hat\varphi=1.11$ in the form $\sigma w+\varphi\Lambda=1$
    (Fig.~\ref{fig:res-planes-social}b). The grid brackets the boundary coarsely: at $w=0$ colonies go extinct for
    $r\le0.05$ (0.9 at $r=0.07$) and survive for $r\ge0.10$; at $w=0.1$ the change happens between $r=0.03$ and 0.05;
    at $w\ge0.2$ no colony goes extinct. In the subcritical corner the median peak is 6\,500--11\,500 individuals at
    weeks 77--96 ($\rho_{\rm pk}$ between 0.85 and 1.66 over the 120 runs), and O4 fails in every run.
\end{itemize}

\begin{table}[ht]
  \centering
  \caption{The confirmatory family with the primary evaluator (left) and with the preregistered evaluator (right), on
  the same runs. Holm correction over seven tests within each column set; the $p$ of Q6 is quasi-binomial. The
  preregistered evaluator takes the transplant of Q1 at $1.75\,t_{\rm pk}$ with $t_{\rm pk}=\arg\max N$, as
  preregistered.}
  \label{tab:res-Qcmp}
  \scriptsize
  \resizebox{\textwidth}{!}{%
  \begin{tabular}{@{}lp{5.2cm}lccp{5.2cm}lcc@{}}
    \toprule
    & \multicolumn{4}{c}{primary evaluator} & \multicolumn{4}{c}{preregistered evaluator} \\
    \cmidrule(lr){2-5}\cmidrule(lr){6-9}
    & Estimate & $p_{\rm Holm}$ & Rej. & Met & Estimate & $p_{\rm Holm}$ & Rej. & Met \\
    \midrule
    Q1 & 0/60 vs.\ 51/60 ($P_T=0.85$ [0.74, 0.92]); discordant pairs 51:0 & $1.3\times10^{-15}$ & yes & yes & 0/60 vs.\ 12/60 ($P_T=0.20$ [0.12, 0.32]); discordant pairs 12:0 & $7.3\times10^{-4}$ & yes & no \\
Q2 & 0.99 vs.\ $-$R1 0.21, $-$R2 0.00, $-$AV 0.79, $-$R7 0.00 & $2.2\times10^{-7}$ & yes & no & 0.97 vs.\ $-$R1 0.23, $-$R2 0.17, $-$AV 0.90, $-$R7 0.00 & 0.067 & no & no \\
Q3 & 200/200 vs.\ 47/100; CV 0.066 [0.059, 0.071] & $<10^{-16}$ & yes & yes & 200/200 vs.\ 47/100; CV 0.066 [0.059, 0.071] & $<10^{-16}$ & yes & yes \\
Q4 & slope $0.996\pm0.014$; $\bar r=-2.0\pm0.3$ wk & $<10^{-16}$ & yes & yes & slope $0.996\pm0.014$; $\bar r=-2.0\pm0.3$ wk & $<10^{-16}$ & yes & yes \\
Q5 & $b_2=+185.3\pm64.3$ (convex); no interior maximum & 1.00 & no & no & $b_2=+69.5\pm14.2$ (convex); no interior maximum & 1.00 & no & no \\
Q6 & $\hat\theta=-0.15$ [-0.20, -0.05] (quasi-binomial, $\phi=3.1$) & $<10^{-16}$ & yes & no & $\hat\theta=0.30$ [0.15, 0.50] (quasi-binomial, $\phi=1.8$) & $4.0\times10^{-5}$ & yes & yes \\
Q7 & 120/120 vs.\ 0/30 (30/30 explode) & $<10^{-16}$ & yes & yes & 120/120 vs.\ 0/30 (30/30 explode) & $<10^{-16}$ & yes & yes \\
\bottomrule

  \end{tabular}}
\end{table}

\paragraph{Q7 with the window varied.}
Fig.~\ref{fig:res-q7aw} shows $P_{\rm ext}$ at $\gamma=0$ against $\Lambda=r\,a_w$ for $a_w\in\{6,12,24\}$ and
$w\in\{0,0.05,0.1\}$: six values of $\Lambda$ per window (0.24, 0.36, 0.48, 0.6, 0.84, 1.2; consecutive ratios
1.25--1.5), 54 points, 20 runs each. The analysis fixed before the production runs is a logistic model in $z=\log r+\theta\log a_w$ and $w$,
with a profile-likelihood interval for $\theta$, a likelihood ratio for $\theta=1$ and an equivalence test of the
two exponents (TOST, margin $[2/3,3/2]$). The data do not support that analysis: 51 of the 54 points have
$P_{\rm ext}=0$ or 1 (the other three are 19/20, 14/20 and 5/20), the Pearson dispersion is 0.05 and the maximum
likelihood coefficients diverge, that is, the model is separated. In a separated logistic model the profile
$\ell(\theta)$ is flat wherever $\theta$ keeps the ordering of the points in $z$ that separates zeros from ones and
drops where it breaks it, so the profile interval ($\hat\theta=0.725$ [0.70, 0.75], likelihood ratio 219 for
$\theta=1$) reflects the geometry of the grid rather than sampling error; the fitted
transition widths are $10^{-8}$--$10^{-19}$ in $\log\Lambda$, below the grid spacing at every $a_w$, and the parametric
bootstrap of the fitted boundary ($\hat\sigma=4.3$, $\hat\varphi=1.17$) resamples essentially the same data. None of
these quantities is reported as inference; they are kept, labelled, in \texttt{results\_numbers.json}.
Table~\ref{tab:res-q7aw} gives instead a non-parametric summary per window: the last grid point with
$P_{\rm ext}\ge1/2$ and the first with $P_{\rm ext}<1/2$, which bracket $\Lambda_{1/2}$ to a factor 1.25--1.5; a
log-linear interpolation between them; and, between consecutive windows, the exponent
$\theta_{\rm loc}=1-\log[\Lambda_{1/2}(a_2)/\Lambda_{1/2}(a_1)]/\log(a_2/a_1)$ implied by the brackets (its grid
bounds) and by the interpolated values. At every $w$ the brackets of $a_w=12$ and 24 coincide and that of $a_w=6$ is
one grid step below, so $\theta_{\rm loc}\in(0,1)$ between 6 and 12 weeks (interpolated 0.51, 0.60 and 0.70 for
$w=0$, 0.05 and 0.1) while between 12 and 24 weeks the grid only bounds it within about $(0.5,1.5)$ (interpolated
0.99, 1.00 and 0.88). The equivalence margin $[2/3,3/2]$ is thus resolved in neither interval, and the one firm
statement is qualitative: at a given product $r\,a_w$ a short window with fast learning is less effective than a
long window with slow learning, and a single exponent does not describe the three windows. Resolving $\theta$
would need a grid of ratio about 1.05 around each $\Lambda_{1/2}$ with more runs per point; it was not run.

\begin{figure}[ht]
  \centering
  \includegraphics[width=\textwidth]{fig_res_q7aw.pdf}
  \caption{Extinction probability without crowding damage ($\gamma=0$) against $\Lambda=r\,a_w$ for three
  plasticity windows, at three inheritance weights (95\% Wilson intervals, 20 runs per point).}
  \label{fig:res-q7aw}
\end{figure}

\begin{table}[ht]
  \centering
  \caption{Transition of $P_{\rm ext}$ at $\gamma=0$ in $\Lambda=r\,a_w$, by inheritance weight $w$ and window
  $a_w$ (20 runs per point). Columns: the last grid value of $\Lambda$ with $P_{\rm ext}\ge1/2$ and the first with
  $P_{\rm ext}<1/2$, each with its count of extinct runs, which bracket $\Lambda_{1/2}$; its log-linear interpolation;
  and the exponent $\theta_{\rm loc}=1-\log[\Lambda_{1/2}(a')/\Lambda_{1/2}(a_w)]/\log(a'/a_w)$ between this window
  and the next ($a'=12$ for $a_w=6$, $a'=24$ for $a_w=12$): the range allowed by the two brackets, then the value
  from the interpolated $\Lambda_{1/2}$. The preregistered logistic fit is separated (51 of 54 points at
  $P_{\rm ext}\in\{0,1\}$) and is not reported as inference (text).}
  \label{tab:res-q7aw}
  \small
  \begin{tabular}{@{}cclllc@{}}
    \toprule
    $w$ & $a_w$ & last $\Lambda$, $P_{\rm ext}\ge\frac12$ & first $\Lambda$, $P_{\rm ext}<\frac12$ & $\Lambda_{1/2}$ & $\theta_{\rm loc}$ to next window \\
    \midrule
    0 & 6 & 0.6 (20/20) & 0.84 (0/20) & 0.71 & (0.00, 1.00); 0.51 \\
0 & 12 & 0.84 (19/20) & 1.2 (0/20) & 0.99 & (0.49, 1.51); 0.99 \\
0 & 24 & 0.84 (20/20) & 1.2 (0/20) & 1.00 & -- \\
\midrule
0.05 & 6 & 0.48 (20/20) & 0.6 (0/20) & 0.54 & (0.19, 1.00); 0.60 \\
0.05 & 12 & 0.6 (20/20) & 0.84 (0/20) & 0.71 & (0.51, 1.49); 1.00 \\
0.05 & 24 & 0.6 (20/20) & 0.84 (0/20) & 0.71 & -- \\
\midrule
0.1 & 6 & 0.36 (20/20) & 0.48 (0/20) & 0.42 & (0.26, 1.00); 0.70 \\
0.1 & 12 & 0.48 (14/20) & 0.6 (0/20) & 0.51 & (0.68, 1.32); 0.88 \\
0.1 & 24 & 0.48 (20/20) & 0.6 (5/20) & 0.56 & -- \\
\bottomrule

  \end{tabular}
\end{table}

\subsection{Re-scoring with the preregistered evaluator and with each option changed separately}
\label{app:res-rescoring}

Every run of the production program was scored with the primary evaluator, with the full preregistered evaluator and with
each option changed separately: peak time at $\arg\max N$ or at the end of the plateau instead of its start; crowded
class at $m>m_c$ instead of $m>8$; O5 on the last sustained birth instead of $t_{B99}$; the preregistered O4 (no
free-space condition). Table~\ref{tab:res-rescoring} gives O17p for the main arms, and Table~\ref{tab:res-null-v1}
the discriminating power of the redefined criteria with the preregistered evaluator.
\begin{itemize}
  \item The peak-time definition matters only through O5. With $t'_{\rm pk}$ and the last sustained birth, O5 fails
    in half of the \textsf{C1} runs (O17p 0.48): the thin tail of the last births lies more than $0.5\,t'_{\rm pk}$
    after the start of the plateau. Taking the end of effective natality at 99\% of the births removes this
    artefact (0.995). With $\arg\max N$ the result is insensitive to the birth reference (1.00). The choice of
    $t_{B99}$ was made after exploratory runs with the primary evaluator (smoke tests on the data of the preregistered design and on two
    replicates per point) had shown this failure of the alternative, so it is informed by the result, and it is the
    one option on which the rejection of Q2 depends (Table~\ref{tab:deviations}).
  \item The fixed crowded-class threshold ($m>8$) changes nothing at the \textsf{C1} values, where $m_c=8$; it
    lowers \textsf{C2} from 0.78 to 0.70 together with the primary O4.
  \item The primary O4 (free space) fails only without attraction and in parts of the $(\alpha,\kappa)$ plane and of
    the hypercubes.
  \item The primary O10 and O11 remove the vacuous passes of the preregistered criteria: the preregistered O10
    passed in every null model, and the preregistered O11 in 1.00 of the $\gamma=0$ runs and in 0.09--0.21 of the
    other null models. The primary O11 also turns the preregistered value of $-$R2 (0.17) into 0 and that of $s_{\rm base}=0.25$ (0.64) into
    0.08.
  \item For the demographic variants the drop in O17p is the same with both evaluators ($\mu_{\rm bg}=0.003$: 0.55
    against 0.45); with $\arg\max N$ as peak time it is masked (0.88), because $t_{\rm pk}$ then moves to the end of
    natality.
\end{itemize}

\begin{table}[ht]
  \centering
  \caption{O17p by arm and evaluator (fraction of runs). Primary: $t'_{\rm pk}$, $t_{B99}$, O4 with free space,
  $m>8$, O10 on adults, O11 on the second half of the births. Prereg.: the full preregistered evaluator. The other columns
  change one option of the primary evaluator: peak time at $\arg\max N$ or at the end of the plateau; crowded class
  at $m>m_c$; O5 on the last sustained birth; the preregistered O4.}
  \label{tab:res-rescoring}
  \scriptsize
  \setlength{\tabcolsep}{4pt}
  \begin{tabular}{@{}lccccccc@{}}
    \toprule
    Arm & primary & prereg. & $\arg\max$ & plateau end & $m>m_c$ & O5 on $t_{\rm last}$ & O4 prereg. \\
    \midrule
    C1 & 0.99 & 0.98 & 1.00 & 1.00 & 0.99 & 0.48 & 0.99 \\
C1$'$ & 0.83 & 0.93 & 0.96 & 0.96 & 0.83 & 0.36 & 0.83 \\
C2 & 0.70 & 0.88 & 0.73 & 0.81 & 0.70 & 0.23 & 0.78 \\
C1 $-$ AV & 0.79 & 0.90 & 0.96 & 0.97 & 0.79 & 0.44 & 0.79 \\
C1 $-$ R1 & 0.21 & 0.23 & 0.22 & 0.22 & 0.21 & 0.12 & 0.21 \\
C1 $-$ R2 & 0.00 & 0.17 & 0.00 & 0.00 & 0.00 & 0.00 & 0.00 \\
C1 $-$ R7 & 0.00 & 0.00 & 0.00 & 0.00 & 0.00 & 0.00 & 0.00 \\
\midrule
C0 & 0.00 & 0.00 & 0.00 & 0.00 & 0.00 & 0.00 & 0.00 \\
C1, $\alpha=0$ & 0.00 & 0.01 & 0.00 & 0.00 & 0.00 & 0.00 & 0.00 \\
$\gamma=0$ & 0.00 & 0.00 & 0.00 & 0.00 & 0.00 & 0.00 & 0.00 \\
$-$R2$^*$ ($w=1$) & 0.00 & 0.15 & 0.00 & 0.00 & 0.00 & 0.00 & 0.00 \\
long-lived & 0.00 & 0.00 & 0.00 & 0.00 & 0.00 & 0.00 & 0.00 \\
no rules & 0.00 & 0.00 & 0.00 & 0.00 & 0.00 & 0.00 & 0.00 \\
\midrule
C1, single-pass R7 & 0.99 & 0.97 & 1.00 & 1.00 & 0.99 & 0.51 & 0.99 \\
C1, $\mu_{\rm bg}=0.003$ & 0.55 & 0.45 & 0.88 & 0.99 & 0.55 & 0.04 & 0.55 \\
C1, $\mu_{\rm bg}=0.008$ & 0.06 & 0.02 & 0.32 & 0.72 & 0.06 & 0.00 & 0.06 \\
C1, nest 3 wk & 0.83 & 0.87 & 0.90 & 0.90 & 0.83 & 0.47 & 0.83 \\
C1, asynchronous & 1.00 & 1.00 & 1.00 & 1.00 & 1.00 & 0.71 & 1.00 \\
C1, async., soft cap. & 0.99 & 0.99 & 0.99 & 0.99 & 0.99 & 0.72 & 0.99 \\
C1, $s_{\rm base}=0.25$ & 0.08 & 0.64 & 0.07 & 0.08 & 0.08 & 0.04 & 0.08 \\
C1, $s_{\rm base}=0.5$ & 0.00 & 0.36 & 0.00 & 0.00 & 0.00 & 0.00 & 0.00 \\
C1, ages U\{4..52\} & 0.93 & 0.90 & 1.00 & 1.00 & 0.93 & 0.36 & 0.93 \\
C1, ages U\{4..80\} & 0.63 & 0.56 & 0.97 & 1.00 & 0.63 & 0.10 & 0.63 \\
\midrule
C1$'$, $\mu_{\rm bg}=0.003$ & 0.28 & 0.28 & 0.68 & 0.89 & 0.28 & 0.02 & 0.28 \\
C1$'$, nest 3 wk & 0.60 & 0.63 & 0.69 & 0.69 & 0.60 & 0.34 & 0.60 \\
C1$'$, asynchronous & 0.94 & 0.98 & 0.99 & 0.99 & 0.94 & 0.39 & 0.94 \\
C1$'$, $s_{\rm base}=0.25$ & 0.26 & 0.53 & 0.30 & 0.30 & 0.26 & 0.17 & 0.26 \\
C1$'$, ages U\{4..52\} & 0.68 & 0.70 & 0.97 & 0.99 & 0.68 & 0.23 & 0.68 \\
C1$'$, ages U\{4..80\} & 0.31 & 0.30 & 0.77 & 0.94 & 0.31 & 0.03 & 0.31 \\
\bottomrule

  \end{tabular}
\end{table}

\begin{table}[ht]
\centering
\caption{Discriminating power of all criteria, essential (top, as in Table~\ref{tab:res-null} of the main text) and
desirable (bottom); primary evaluator, fraction of all runs that pass, NA counting as failure. Columns: the
reference ensemble (\textsf{C1}, \textsf{C1}$'$, $\alpha=0$, \textsf{C0}), the four null models with the initial
condition of \textsf{C1} ($\gamma=0$; $-$R2$^*$: $w=1$ without social learning; long-lived control; no rules: \textsf{C1} without its four
rules), $n=200$ each, and the short-lived control ($n=100$). ``Generic'': passes in at least two of the four null models.
``Discriminates'': upper Wilson bound below 0.8 and Newcombe interval of the difference with \textsf{C1} excluding
zero (long: long-lived control; none: no rules; short: short-lived control).}
\label{tab:res-null-full}
\footnotesize
\setlength{\tabcolsep}{3.2pt}
\begin{tabular}{@{}lccccccccccl@{}}
\toprule
& \textsf{C1} & \textsf{C1}$'$ & $\gamma{=}0$ & $\alpha{=}0$ & \textsf{C0} & $-$R2$^*$ & long & none & short & generic & discriminates against \\
\midrule
O4 & 1.00 & 1.00 & 0.00 & 0.00 & 1.00 & 0.99 & 0.17 & 1.00 & 0.87 & yes & $\gamma{=}0$, $\alpha{=}0$, long \\
O5 & 0.99 & 0.85 & 1.00 & 1.00 & 1.00 & 0.97 & 0.89 & 0.91 & 0.25 & yes & short \\
O6 & 1.00 & 1.00 & 0.00 & 1.00 & 1.00 & 1.00 & 0.99 & 1.00 & 0.45 & yes & $\gamma{=}0$, short \\
O7 & 1.00 & 0.99 & 0.00 & 1.00 & 1.00 & 1.00 & 0.99 & 1.00 & 0.47 & yes & $\gamma{=}0$, short \\
O10 & 1.00 & 1.00 & 0.00 & 1.00 & 1.00 & 1.00 & 1.00 & 1.00 & 1.00 & yes & $\gamma{=}0$ \\
O11 & 1.00 & 0.97 & 0.00 & 1.00 & 1.00 & 0.00 & 0.00 & 0.00 & 0.00 & no & $\gamma{=}0$, $-$R2, long, none, short \\
O13 & 1.00 & 1.00 & 0.00 & 1.00 & 0.92 & 1.00 & 0.00 & 0.00 & 0.00 & no & $\gamma{=}0$, long, none, short \\
O13p & 1.00 & 1.00 & 0.00 & 0.00 & 0.00 & 1.00 & 0.00 & 0.00 & 0.00 & no & $\gamma{=}0$, $\alpha{=}0$, C0, long, none, short \\
O17p & 0.99 & 0.83 & 0.00 & 0.00 & 0.00 & 0.00 & 0.00 & 0.00 & 0.00 & no & $\gamma{=}0$, $\alpha{=}0$, C0, $-$R2, long, none, short \\
\midrule
O10b & 0.02 & 0.01 & 1.00 & 1.00 & 0.18 & 0.05 & 0.20 & 0.12 & 1.00 & no & none \\
O8 & 0.62 & 0.21 & 0.00 & 0.97 & 0.82 & 0.69 & 0.39 & 0.64 & 0.04 & no & C1$'$, $\gamma{=}0$, long, short \\
O8b & 1.00 & 1.00 & 0.00 & 1.00 & 1.00 & 1.00 & 1.00 & 1.00 & 1.00 & yes & $\gamma{=}0$ \\
O9 & 0.45 & 0.29 & 0.00 & 0.49 & 0.44 & 0.51 & 0.39 & 0.45 & 0.18 & no & C1$'$, $\gamma{=}0$, short \\
O11b & 1.00 & 1.00 & 1.00 & 1.00 & 1.00 & 0.99 & 0.95 & 0.93 & 1.00 & yes & none \\
O14 & 1.00 & 1.00 & 1.00 & 0.00 & 1.00 & 1.00 & 1.00 & 1.00 & 1.00 & yes & $\alpha{=}0$ \\
\bottomrule

\end{tabular}
\end{table}

\begin{table}[ht]
  \centering
  \caption{Discriminating power of the redefined criteria with the preregistered evaluator (same ensembles and columns
  as Table~\ref{tab:res-null}). The preregistered O10 passes in every model; the preregistered O11 passes vacuously without
  crowding damage.}
  \label{tab:res-null-v1}
  \footnotesize
  \setlength{\tabcolsep}{3.2pt}
  \begin{tabular}{@{}lccccccccccl@{}}
    \toprule
    & \textsf{C1} & \textsf{C1}$'$ & $\gamma{=}0$ & $\alpha{=}0$ & \textsf{C0} & $-$R2$^*$ & long & none & short & generic & discriminates against \\
    \midrule
    O4 & 1.00 & 1.00 & 0.00 & 1.00 & 1.00 & 1.00 & 0.17 & 1.00 & 0.87 & yes & $\gamma{=}0$, long \\
O5 & 1.00 & 0.99 & 0.00 & 1.00 & 1.00 & 1.00 & 0.96 & 0.99 & 0.16 & yes & $\gamma{=}0$, short \\
O10 & 1.00 & 1.00 & 1.00 & 1.00 & 1.00 & 1.00 & 1.00 & 1.00 & 1.00 & yes & none \\
O11 & 0.98 & 0.94 & 1.00 & 1.00 & 0.98 & 0.15 & 0.09 & 0.21 & 0.00 & no & $-$R2, long, none, short \\
O13p & 1.00 & 1.00 & 0.00 & 0.01 & 0.00 & 1.00 & 0.00 & 0.00 & 0.00 & no & $\gamma{=}0$, $\alpha{=}0$, C0, long, none, short \\
O17p & 0.98 & 0.93 & 0.00 & 0.01 & 0.00 & 0.15 & 0.00 & 0.00 & 0.00 & no & $\gamma{=}0$, $\alpha{=}0$, C0, $-$R2, long, none, short \\
\bottomrule

  \end{tabular}
\end{table}

\paragraph{Threshold of the deteriorated classes.}
\label{app:res-threshold}
O13 requires both the crowded and the isolated deteriorated classes to hold at least $\theta=10\%$ of the adults,
and O13p the same of the persistent isolated class. The stored class fractions at $t'_{\rm pk}$, $t'_{\rm pk}+10$
and $t'_{\rm pk}+20$ allow both to be re-scored for other $\theta$ without new simulations; the other patterns of
O17p are left unchanged. Fig.~\ref{fig:res-theta} gives O17p against $\theta$ for the candidates, two variants and
the single ablations, on the grid $\theta\in\{0.05,\dots,0.20\}$ against the ensemble thresholds
$\tau\in\{0.5,\dots,0.9\}$. \textsf{C1} holds 0.99 up to $\theta=0.125$ and 0.96 at 0.15, and
collapses at 0.175; \textsf{C1}$'$ holds 0.83 up to 0.125 and falls to 0.33 at 0.15 (0.31 in the $-$AV arm of the
factorial, $n=120$); \textsf{C2} falls from 0.70 to
0.01 at 0.125; the asynchronous variant and the nest variant of \textsf{C1} hold up to 0.125 and 0.15. R1, R2 and R7
are necessary at every $(\theta,\tau)$ of the grid; AV is necessary for $\tau\ge0.8$ when $\theta\le0.125$ (0.79
against 0.8, with a Wilson interval [0.71, 0.85] that contains the threshold) and for every $\tau$ when
$\theta\ge0.15$. In the hypercubes, the Calhoun-like fraction is 5.5\%, 2.3\% and 0.5\% at $\theta=0.05$, 0.10 and
0.15. The conclusions about refuges are therefore conditional on the 10\% threshold, through the capacity bound
$\Phi^{\rm pers}_{\rm BO}\le\phi_R$.

\begin{figure}[ht]
  \centering
  \includegraphics[width=0.8\textwidth]{fig_res_theta.pdf}
  \caption{O17p re-scored with the class threshold $\theta$ (the production value is 0.10, dotted).
  \textbf{(a)}~Candidates and two variants of \textsf{C1}. \textbf{(b)}~\textsf{C1} and its single ablations
  ($n=120$); the horizontal lines are the ensemble thresholds $\tau=0.5$--0.9 of the necessity grid.}
  \label{fig:res-theta}
\end{figure}

\subsection{Robustness variants and system size: tables}
\label{app:res-rob}

Table~\ref{tab:res-rob} lists the robustness variants of Section~\ref{sec:res-robustness} and
Table~\ref{tab:res-size} the lattice sizes at fixed density. In every variant that loses the primary outcome the
loss comes from a single criterion: O5 for the background mortality and the non-synchronous initial conditions
(and, marginally, for \textsf{C1}$'$), O11 for the nest period and for pups born above the deterioration threshold.
O5 is lost through its timing clause $(t_{B99}-t'_{\rm pk})/t'_{\rm pk}\le0.5$, not through its level clause
$N(t_{B99})\ge0.7\,N_{\rm pk}$: at $\mu_{\rm bg}=0.003$ the O5 verdicts of \textsf{C1} are 110 PASS, 83 MARGINAL and
7 FAIL (median $t'_{\rm pk}$ 29 weeks against 34 in the reference ensemble, median $t_{B99}$ 43 against 40); at
0.008, 12, 118 and 70 ($t'_{\rm pk}$ 26, $t_{B99}$ 50); with ages U\{4..80\}, 126, 73 and 1 ($t'_{\rm pk}$ 29,
$t_{B99}$ 41); with U\{4..52\}, 186, 14 and 0. The column ``disc.'' of Table~\ref{tab:res-rob} shows that the
criteria that discriminate against the null models (with O11 scored at six weeks) are untouched by the demographic variants (1.00 in every
\textsf{C1} arm with background mortality or spread ages; 0.96--0.99 in \textsf{C1}$'$) and are lost only where O11 is
lost (nest period, $s_{\rm base}>0$). Along the background-mortality curve the primary outcome is lost gradually: at
$\mu_{\rm bg}=0.001$ per week, a tenth of the senescent hazard, O17p is still 0.83. In the lattice-size stage the
failures at $L=15$ (11 of 60) and $L=20$ (6 of 60) are O5 marginal in 5 runs at each size, O11 marginal in 1 at
each, and at $L=15$ the anti-pattern A2 (recovery of natality) in 6 runs; at $L=30$ one run fails O5 and at
$L\ge45$ none fails.

\begin{figure}[ht]
  \centering
  \includegraphics[width=\textwidth]{fig_res_rob.pdf}
  \caption{Robustness of the primary outcome ($n=200$ per arm unless stated).
  \textbf{(a)}~O17p against the background mortality $\mu_{\rm bg}$ for \textsf{C1} and \textsf{C1}$'$ ($n=60$ at the
  intermediate values), with the primary evaluator (filled), the preregistered evaluator (open) and with $\arg\max N$
  as peak time (crosses).
  \textbf{(b)}~Variants of \textsf{C1} and \textsf{C1}$'$: single-pass refuge admission, background mortality, nest
  period, asynchronous update with hard and soft capacity, pups born at $s_{\rm base}=0.25$ and 0.5, and initial
  ages drawn from U\{4..52\} and U\{4..80\}. Vertical lines: the reference ensembles and the 0.8 threshold.
  \textbf{(c)}~Lattice side at fixed density ($n=60$): KM spread of the extinction time, relative spread of the
  last birth and O17p.}
  \label{fig:res-rob}
\end{figure}

\begin{table}[ht]
  \centering
  \caption{Robustness variants ($n=200$ per arm). O17p with 95\% Wilson interval and Newcombe interval of the
  difference with the reference ensemble; the same arm with the preregistered evaluator and with $\arg\max N$ as peak
  time; ``disc.'': O17p restricted to the criteria that discriminate against the null models (O11 scored at six
  weeks, O13 and O13p pass and no anti-pattern is triggered), which leaves out the generic O4--O7 and O10; the criterion with the lowest
  pass rate when O17p $<0.95$ (NA counted as failure); median peak density, plateau length (weeks within 2\% of the
  maximum) and Kaplan--Meier median of $T_{\rm ext}$.}
  \label{tab:res-rob}
  \scriptsize
  \setlength{\tabcolsep}{3.2pt}
  \begin{tabular}{@{}lllccclccc@{}}
    \toprule
    Variant & O17p & $\Delta$ vs.\ reference & prereg. & $\arg\max$ & disc. & fails & $\rho_{\rm pk}$ & plateau & $T_{\rm ext}$ \\
    \midrule
    \emph{C1}, reference & 0.99 [0.97, 1.00] & -- & 0.98 & 1.00 & 1.00 & -- & 0.36 & 51 & 169 \\
single-pass R7 admission & 0.99 [0.97, 1.00] & +0.00 [-0.02, +0.02] & 0.97 & 1.00 & 1.00 & -- & 0.37 & 52 & 168 \\
$\mu_{\rm bg}=0.003$ & 0.55 [0.48, 0.62] & -0.44 [-0.51, -0.37] & 0.45 & 0.88 & 1.00 & O5 0.55 & 0.33 & 17 & 174 \\
$\mu_{\rm bg}=0.008$ & 0.06 [0.03, 0.10] & -0.94 [-0.96, -0.89] & 0.02 & 0.32 & 1.00 & O5 0.06 & 0.30 & 10 & 182 \\
nest period 3 wk & 0.83 [0.78, 0.88] & -0.16 [-0.22, -0.11] & 0.87 & 0.90 & 0.90 & O11 0.90 & 0.31 & 53 & 171 \\
asynchronous, hard & 1.00 [0.98, 1.00] & +0.01 [-0.01, +0.03] & 1.00 & 1.00 & 1.00 & -- & 0.43 & 57 & 159 \\
asynchronous, soft & 0.99 [0.97, 1.00] & +0.00 [-0.02, +0.02] & 0.99 & 0.99 & 0.99 & -- & 0.42 & 57 & 158 \\
$s_{\rm base}=0.25$ & 0.08 [0.05, 0.13] & -0.92 [-0.95, -0.86] & 0.64 & 0.07 & 0.08 & O11 0.08 & 0.40 & 53 & 164 \\
$s_{\rm base}=0.5$ & 0.00 [0.00, 0.02] & -0.99 [-1.00, -0.97] & 0.36 & 0.00 & 0.00 & O11 0.00 & 0.44 & 53 & 165 \\
ages U\{4..52\} & 0.93 [0.89, 0.96] & -0.06 [-0.11, -0.03] & 0.90 & 1.00 & 1.00 & O5 0.93 & 0.36 & 28 & 167 \\
ages U\{4..80\} & 0.63 [0.56, 0.69] & -0.36 [-0.43, -0.30] & 0.56 & 0.97 & 1.00 & O5 0.63 & 0.34 & 22 & 173 \\
\midrule
\emph{C1$'$}, reference & 0.83 [0.78, 0.88] & -- & 0.93 & 0.96 & 0.97 & -- & 0.32 & 50 & 182 \\
single-pass R7 admission & 0.85 [0.79, 0.89] & +0.02 [-0.06, +0.09] & 0.91 & 0.95 & 0.96 & O5 0.89 & 0.33 & 49 & 180 \\
$\mu_{\rm bg}=0.003$ & 0.28 [0.22, 0.34] & -0.56 [-0.63, -0.47] & 0.28 & 0.68 & 0.96 & O5 0.28 & 0.29 & 19 & 183 \\
$\mu_{\rm bg}=0.008$ & 0.04 [0.02, 0.08] & -0.79 [-0.84, -0.73] & 0.01 & 0.24 & 0.96 & O5 0.04 & 0.26 & 9 & 182 \\
nest period 3 wk & 0.60 [0.54, 0.67] & -0.23 [-0.31, -0.14] & 0.63 & 0.69 & 0.69 & O11 0.69 & 0.30 & 52 & 181 \\
asynchronous, hard & 0.94 [0.89, 0.96] & +0.10 [+0.04, +0.16] & 0.98 & 0.99 & 0.99 & O5 0.94 & 0.38 & 56 & 173 \\
asynchronous, soft & 0.97 [0.94, 0.99] & +0.14 [+0.08, +0.19] & 0.98 & 1.00 & 1.00 & -- & 0.38 & 56 & 173 \\
$s_{\rm base}=0.25$ & 0.26 [0.20, 0.32] & -0.57 [-0.65, -0.49] & 0.53 & 0.30 & 0.31 & O11 0.32 & 0.38 & 50 & 175 \\
$s_{\rm base}=0.5$ & 0.00 [0.00, 0.02] & -0.83 [-0.88, -0.77] & 0.38 & 0.01 & 0.00 & O11 0.01 & 0.43 & 47 & 174 \\
ages U\{4..52\} & 0.68 [0.61, 0.74] & -0.15 [-0.24, -0.07] & 0.70 & 0.97 & 0.99 & O5 0.69 & 0.33 & 26 & 182 \\
ages U\{4..80\} & 0.31 [0.25, 0.38] & -0.52 [-0.60, -0.44] & 0.30 & 0.77 & 0.99 & O5 0.31 & 0.30 & 21 & 187 \\
\bottomrule

  \end{tabular}
\end{table}

\begin{table}[ht]
  \centering
  \caption{Lattice side $L$ at fixed seeding density $N_0/L^2=0.44$ ($n=60$ per size, $T=600$). Medians over runs;
  $t_{\rm last}$ is the last surviving birth, $D_{\rm last}=T_{\rm ext}-t_{\rm last}$; the spread is the
  interquartile range over the median of the Kaplan--Meier estimate of $T_{\rm ext}$, with bootstrap interval.}
  \label{tab:res-size}
  \scriptsize
  \setlength{\tabcolsep}{3.5pt}
  \begin{tabular}{@{}rrrlcccrrrrl@{}}
    \toprule
    $L$ & $N_0$ & $n$ & O17p & $\rho_{\rm pk}$ & $f_0$ & $\Phi^{\rm pers}_{\rm BO}$ & $t_{\rm last}$ & IQR & $D_{\rm last}$ & $T_{\rm ext}$ & spread \\
    \midrule
    15 & 100 & 60 & 0.82 [0.70, 0.89] & 0.35 & 0.18 & 0.18 & 50 & 29 & 110 & 162 & 0.148 [0.121, 0.194] \\
20 & 178 & 60 & 0.90 [0.80, 0.95] & 0.36 & 0.17 & 0.17 & 62 & 22 & 112 & 175 & 0.143 [0.110, 0.169] \\
30 & 400 & 60 & 0.98 [0.91, 1.00] & 0.37 & 0.16 & 0.17 & 58 & 13 & 110 & 170 & 0.076 [0.053, 0.112] \\
45 & 900 & 60 & 1.00 [0.94, 1.00] & 0.37 & 0.15 & 0.17 & 66 & 16 & 111 & 176 & 0.085 [0.052, 0.113] \\
60 & 1600 & 60 & 1.00 [0.94, 1.00] & 0.37 & 0.15 & 0.17 & 75 & 14 & 110 & 184 & 0.076 [0.050, 0.103] \\
\bottomrule

  \end{tabular}
\end{table}

\subsection{Phase planes}
\label{sec:res-planes}

All planes use 30 runs per point; the 40 points nearest to a Calhoun-like boundary (9 of P1, 12 of P1-$a_w$ and 19
of P2) were re-run with 100 independent runs.

\paragraph{Socialisation: inheritance $w$ and learning rate $r$ (P1).}
With crowding damage ($\gamma=0.1$, Fig.~\ref{fig:res-planes-social}a), every one of the 56 points goes extinct
($P_{\rm ext}=1$), so the plane decides only \emph{which} collapse occurs. O17p is 0.83--1.00 for $w\le0.3$ and
$r\le0.10$. It falls in two directions, both through the primary O11: with strong inheritance ($w\ge0.5$: 0.33 at
$r=0.02$, 0 at $w\ge0.75$) the late cohorts inherit the competence of their mothers, and with fast learning
($r\ge0.2$: 0.57 at $w=0.2$, $r=0.2$; 0 at $r=0.3$ for $w\ge0.2$) they acquire it from the few competent adults.
These are point estimates with 30 runs per cell (Wilson half-width up to $\pm0.17$); the boundaries of the maps have
no intervals of their own. Without crowding damage ($\gamma=0$, panel b), colonies either explode or, in the
subcritical corner, collapse endogenously after overshooting the capacity (Q7).

\paragraph{Plasticity window $a_w$ and learning rate $r$ (P1-$a_w$).}
For every finite window up to $a_w=24$ weeks and every $r$ from 0.02 to 0.3, extinction is certain (48 of 48
points at $P_{\rm ext}=1$). O17p degrades at the fastest learning rates: at $r=0.2$ it is 0.10--0.13 for
$a_w\le8$ and 0.40--0.60 for $a_w\ge12$, and at $r=0.3$ it is at most 0.07; in every case O11 is the criterion that
fails, more often the shorter the window. The loss of the persistent class with long windows and fast learning,
which set the boundary with the preregistered evaluator (O13p 0.17 at $a_w=24$, $r=0.3$), is hidden behind O11 with the primary evaluator.
Without a window ($a_w=\infty$, top row of Fig.~\ref{fig:res-planes-social}c), recovery is again individual and
reversible: $P_{\rm ext}$ falls from 1.00 at $r=0.02$ to 0.47 at $r=0.07$ and to 0 for $r\ge0.2$, and the
anti-pattern A2 (recovery of natality) rises from 0 to 1.

\paragraph{Attraction $\alpha$ and aversion $\kappa$ (P2).}
O17p is at least 0.8 at 34 of the 48 points with $\alpha\ge2$ and at 35 of the 56 points of the plane
(Fig.~\ref{fig:res-planes-space}a). It fails in three places:
\begin{itemize}
  \item where the free space of the primary O4 is lost: at $\alpha\le2$ with $\kappa\ge0.1$, at $\alpha=3$ with
    $\kappa=0.2$ and at $\kappa=0.3$ for $\alpha\le5$, the empty fraction at the peak is below 0.10 (panel c);
  \item at $\kappa=0.3$ with $\alpha\ge5$, where in addition the crowded fraction falls to the 10\% threshold of O13
    (0.33--0.73);
  \item without aversion and with strong attraction ($\kappa=0$, $\alpha\ge5$; 0.67, 0.63 and 0.33), where O5 and
    O11 fail in a third of the runs; with 100 runs the point $(\alpha,\kappa)=(8,0)$ gives 0.53 [0.43, 0.63].
\end{itemize}
The empty fraction at the peak increases with $\alpha$ and decreases with $\kappa$ (Spearman $+0.56$ and $-0.77$,
exploratory, BH-significant; panel c). The peak density does the opposite. The coexistence band predicted for the
model without refuges is not observed in \textsf{C1} (Q5). Varying the aversion exponent at $\alpha=4$ ($p=2$ or 8)
leaves O17p at or above 0.8 except at $\kappa=0.3$ (0 and 0.30) and at $p=2$, $\kappa=0.2$ (0.47).

\paragraph{Refuges: fraction $f_R$, preference $\pi$ and capacity $c_R$ (P5).}
A persistent isolated class requires both a strong preference and a large enough total capacity: with $c_R=2$,
$\pi\ge30$ and $f_Rc_R\ge0.6$, or $\pi\ge100$ and $f_Rc_R\ge0.4$; with $c_R=1$, $\pi\ge100$ and $f_Rc_R\ge0.4$, or
$\pi=1000$ and $f_Rc_R\ge0.3$ (9 of 50 points with O13p $\ge0.8$; Fig.~\ref{fig:res-planes-space}d,e). The bound
$\Phi^{\rm pers}_{\rm BO}\le\phi_R=f_Rc_R|G|/N_{\rm pk}$ of Section~\ref{sec:model-groups} holds at all 50 points
and is approached for $\pi\ge100$ (panel f). \textsf{C1} lies at $\phi_R=0.21$, with
$\Phi^{\rm pers}_{\rm BO}=0.17$. Large capacity combined with the strongest preference degrades the senescent
decline: O8 is 0.00--0.13 at $\pi=1000$ with $f_R\ge0.2$, $c_R=2$, and O17p falls to 0.47 at $f_R=0.4$ through O5.

\paragraph{Refined boundary points.}
Over the 40 refined points the 100-run estimate differs from the 30-run estimate by 0.07 on average (mean shift
$+0.03$, maximum 0.30, correlation 0.93), the size expected from binomial sampling with 30 runs; 26 of the 40
30-run values lie within the 100-run Wilson interval. Five points crossed the 0.8 threshold, four of them upwards.
The points were selected because their 30-run estimate lay near the boundary, so part of these shifts is
regression to the mean. No conclusion of the planes changes.

\begin{figure}[ht]
  \centering
  \includegraphics[width=\textwidth]{fig_res_planes_social.pdf}
  \caption{Socialisation planes around \textsf{C1} (30 runs per cell, primary evaluator).
  \textbf{(a)}~O17p in the plane of learning rate $r$ and inheritance $w$, with crowding damage $\gamma=0.1$;
  every cell goes extinct.
  \textbf{(b)}~Extinction probability in the same plane without crowding damage ($\gamma=0$); the other colonies
  explode. The dashed line is the fitted boundary $w+\hat\varphi\,r a_w=c$ of the generational map (Q7).
  \textbf{(c)}~O17p in the plane of $r$ and the critical window $a_w$. The top row ($a_w=\infty$) is \textsf{C1}
  without R1.}
  \label{fig:res-planes-social}
\end{figure}

\begin{figure}[ht]
  \centering
  \includegraphics[width=\textwidth]{fig_res_planes_space.pdf}
  \caption{Spatial planes around \textsf{C1} (30 runs per cell, primary evaluator).
  \textbf{(a)--(c)}~Plane of attraction $\alpha$ and aversion $\kappa$ (exponent $p=4$): O17p, the instantaneous
  coexistence criterion O13, and the empty fraction at the peak. The dashed and dotted lines mark $n^*=m_c$ and
  $n^*=4$, with $n^*=1/\kappa-1/\alpha$.
  \textbf{(d, e)}~Persistent O13 in the refuge plane of fraction $f_R$ and preference $\pi$, for capacities
  $c_R=1$ and $2$.
  \textbf{(f)}~Persistent isolated fraction against the refuge capacity per individual at the peak, $\phi_R$, for
  the 50 points of (d, e). The dashed line is the bound $\Phi^{\rm pers}_{\rm BO}=\phi_R$.}
  \label{fig:res-planes-space}
\end{figure}

\subsection{Global sensitivity}
\label{sec:res-sensitivity}

Fig.~\ref{fig:res-sensitivity} reports the elementary effects of the Morris screening, over 15 factors and ranges
around \textsf{C1}, and the Sobol' indices of the 11 factors of the production design, each with its bootstrap interval
and the empirical noise floor of Section~\ref{sec:stats}. Table~\ref{tab:res-region} and Fig.~\ref{fig:res-region}
give the estimators of the Calhoun-like region of the hypercubes.

\paragraph{Noise floor.}
The split-half contrast $e=(\bar Y_{\rm even}-\bar Y_{\rm odd})/2$ of the replicate means at each design point has
expectation zero and the variance of the replicate noise, so the Morris and Sobol' estimators applied to it give the
floor below which an index is indistinguishable from noise. A factor is distinguishable in the screening if its
$\mu^*$ exceeds that of $e$ by a factor of two, and relevant for a Sobol' index if $S_T-S_T^{\rm noise}>0.05$ and the
lower confidence bound exceeds the floor; without the floor, synthetic data with three replicates and null factors
give $S_T\approx0.3$.

\paragraph{Screening.}
For O17p the largest elementary effects are those of the refuge capacity $c_R$ ($\mu^*=0.18$), the aversion
$\kappa$ (0.16), the community size $m_c$ (0.15), the conception probability $p_0$ (0.09) and the attraction
$\alpha$ (0.08); the preference $\pi$ and the inheritance $w$ follow (0.075). With the primary evaluator O17p is
zero at most Morris points, so the elementary effects are small and their bootstrap intervals are as large as the
effects. For every factor $\sigma$ is about twice $\mu^*$; because each output is an average of 10 binary runs,
$\sigma$ includes replicate noise, and without a noise floor for $\sigma$ this does not show that the factors act
through interactions. The automatic selection averages the relative $\mu^*$ over seven metrics and retained $m_c$,
$\kappa$, $\alpha$, $c_R$, $\gamma$, $w$, $\eta$ and $r$, together with the aversion exponent $p$ and $p_0$; $w$, $r$ and
$a_w$ were added by design, so the Sobol' design has 11 factors and holds four at their
\textsf{C1} values: $f_R$, $\pi$, $a_{\max}$ and $\delta$. The indices below are conditional on these four.

\paragraph{Variance decomposition.}
The Saltelli design ($N=256$, 3\,328 points, 11 factors) was run with 3 runs per point and then extended to 6
(19\,968 runs); the values below use the 6 runs. Doubling the runs halved the noise floors (median floor of $S_T$ for
O17p from 0.12 to 0.07) and changed the total indices of the continuous outputs by at most 0.01 ($\rho_{\rm pk}$,
$f_0$, $\Phi^{\rm pers}_{\rm BO}$) and those of O13 and O13p by at most 0.05; for O17p the largest change is 0.10
($m_c$ from 0.39 to 0.31, $\gamma$ from 0.26 to 0.16), and the number of factors above the floor rises from 7 to 10.
With the primary evaluator O17p is zero at 78\% of the points and at least 0.8 at 9\% (mean 0.135): the
Calhoun-like region is a minority of the box even with four parameters held at their \textsf{C1} values. For O17p
the total indices are $c_R$ $0.49\pm0.16$, $w$ $0.43\pm0.15$, $\kappa$ $0.41\pm0.14$, $r$ $0.32\pm0.12$ and $m_c$
$0.31\pm0.12$, followed by $\eta$, $\gamma$, $p$, $p_0$ and $\alpha$ (0.12--0.17, against floors of 0.05--0.09); $a_w$
(0.11) does not exceed its floor. The intervals of the five leading factors overlap, so they cannot be ordered. The
first-order indices add up to 0.15 and the total indices to 2.8: O17p is almost entirely interactions, and 31\% of
its run-level variance is within points. In the Sobol' design of the preregistered plan, which excluded $w$, $r$ and
$a_w$, $m_c$, $f_R$ and $c_R$ led; with the 11 factors, the inheritance $w$ and the aversion $\kappa$ rank with the refuge
capacity, in agreement with the hypercube regressions; with the free-space condition of O4, $\kappa$ enters through
O4 ($S_T=0.55$), and $w$ through O11 ($S_T=0.75$, $S_1=0.49$). For the persistent class O13p the leading factors
are $c_R$ (0.62), $m_c$ (0.48) and $\kappa$ (0.30), and $w$ is below the floor. Two of these are partly
definitional: with $c_R=3$ refuge residents exceed the isolation threshold of O13, so O13p averages 0.56 at $c_R=2$
against 0.11 at $c_R=1$ and 0 at $c_R=3$; and $m_c$ also defines the damage threshold, so the mean of O13 falls from
0.68 for $m_c\in(6,9]$ to 0.39 for $m_c>12$.
The continuous outputs are almost additive and dominated by the spatial and learning parameters: for the peak
density (sum of first-order indices 1.01) $r$ 0.21, $\kappa$ 0.19, $\eta$ 0.15, $\gamma$ 0.13, $a_w$ 0.10 and
$\alpha$ 0.10, every factor above a floor of 0.002; for the empty fraction (0.76) $\kappa$ 0.44, $m_c$ 0.28, $\alpha$
0.16 and $\eta$ 0.12; for $\Phi^{\rm pers}_{\rm BO}$ (0.89) $c_R$ 0.64 and $m_c$ 0.37, with no other factor above the
floor. The indices of O7 and $P_{\rm ext}$ are degenerate (98\% of the points share the same value; total indices
above one with intervals wider than one) and are not reported: over the box, extinction is generic. For the median
extinction time $m_c$ (0.66), $\kappa$ (0.32) and $p_0$ (0.30) exceed floors of about 0.1; with 3 runs per point only
the first two did.

\paragraph{Latin hypercubes.}
A Firth logistic regression of O17p on the standardised factors shows lack of fit (Pearson dispersion 2--3), so
standard errors are inflated by its square root. With this quasi-binomial correction and a false discovery rate of
10\%, the factors significant in the first hypercube are $\kappa$ ($-1.4$), $f_R$ ($+1.4$), $w$ ($-1.4$), $r$
($-1.2$), $\gamma$ ($+1.1$), $\pi$ ($+1.0$), $a_{\max}$ ($+0.8$), $\alpha$ ($+0.7$), $m_c$ ($-0.6$) and $a_w$
($+0.5$); in the second, $\kappa$ ($-1.7$), $f_R$ ($+1.1$), $m_c$ ($-1.0$), $\pi$ ($+0.9$), $p$ ($+0.8$), $w$ ($-0.7$),
$a_{\max}$ ($+0.7$), $\eta$ ($+0.6$), $p_0$ ($+0.5$) and $c_R$ ($-0.4$). The effects of $\kappa$, $f_R$, $\pi$, $w$,
$m_c$ and $a_{\max}$ replicate; those of $r$, $\gamma$, $\alpha$ and $a_w$ appear only in the first hypercube and
those of $p$, $\eta$, $p_0$ and $c_R$ only in the second (the effect of $c_R$ is non-monotone and a linear term does
not describe it). With the preregistered evaluator the leading factors were $m_c$, $f_R$, $\pi$, $w$ and $\kappa$; the
primary O4 and O5 add the aversion and the learning parameters.

\paragraph{Estimators of the Calhoun-like region.}
Table~\ref{tab:res-region} gives, for the primary evaluator, the preregistered evaluator and the fixed class threshold
alone, the naive fraction of points with $k/n\ge0.8$, the mean of O17p over the box and the empirical-Bayes
fraction, without and with the 40 extra runs at the 18 boundary points. The extra runs halve the number of
uncertain points (posterior probability between 0.05 and 0.95) from 12 to 8 and leave the estimates within their
intervals. Taking each option separately, the naive fraction is 2.0\% with the primary evaluator, 3.8\% with
$\arg\max N$ as peak time, 5.0\% with the end of the plateau, 4.2\% with the preregistered O4, 2.5\% with $m>m_c$ and
0\% with O5 on the last sustained birth; the preregistered evaluator gives 7.2\%.

\begin{table}[ht]
  \centering
  \caption{Estimators of the Calhoun-like fraction of the two 15-factor hypercubes (400 points, 10 runs each),
  relative to the box and to the sampling scales. Naive: fraction of points with $k/n\ge0.8$; volume: mean of O17p
  over the box; EB: mean posterior probability of $p_i\ge0.8$ under a beta prior centred on a logistic emulator, with
  dispersion $\phi$ fitted by beta-binomial maximum likelihood. Bootstrap intervals over points with re-fit;
  ``uncertain'': points with posterior probability between 0.05 and 0.95. The extra runs are 40 additional runs at the 18
  points with $0.4\le k/n\le0.95$ under the primary evaluator.}
  \label{tab:res-region}
  \scriptsize
  \setlength{\tabcolsep}{3.5pt}
  \begin{tabular}{@{}lllllcc@{}}
    \toprule
    Evaluator & Runs & naive & volume & EB & $\phi$ & uncertain \\
    \midrule
    primary & 10 per point & 0.022 [0.010, 0.040] & 0.041 [0.027, 0.060] & 0.015 [0.006, 0.029] & 1.22 & 12 \\
primary & +40 at 18 points & 0.020 [0.007, 0.036] & 0.040 [0.026, 0.059] & 0.018 [0.007, 0.032] & 0.82 & 8 \\
preregistered & 10 per point & 0.077 [0.052, 0.105] & 0.110 [0.082, 0.137] & 0.066 [0.045, 0.090] & 0.63 & 30 \\
preregistered & +40 at 18 points & 0.072 [0.050, 0.098] & 0.110 [0.083, 0.137] & 0.069 [0.048, 0.092] & 0.60 & 26 \\
$m>m_c$ & 10 per point & 0.028 [0.013, 0.045] & 0.054 [0.039, 0.075] & 0.020 [0.009, 0.034] & 1.00 & 17 \\
$m>m_c$ & +40 at 18 points & 0.025 [0.013, 0.043] & 0.054 [0.038, 0.073] & 0.021 [0.010, 0.035] & 0.79 & 12 \\
\bottomrule

  \end{tabular}
\end{table}

\begin{figure}[ht]
  \centering
  \includegraphics[width=0.8\textwidth]{fig_res_region.pdf}
  \caption{The Calhoun-like region of the hypercubes with per-point uncertainty (primary evaluator).
  \textbf{(a)}~Posterior probability that a point has $p\ge0.8$ against its observed fraction $k/n$; the 18
  boundary points with 50 runs are highlighted.
  \textbf{(b)}~Fraction of points with $k/n\ge0.8$ in nested subsets chosen after the runs ($c_R\le2$, then
  $m_c\le10$, $f_Rc_R\ge0.3$ and $\pi\ge100$), with bootstrap intervals, for the primary evaluator, the preregistered
  evaluator and the fixed class threshold alone.}
  \label{fig:res-region}
\end{figure}

\begin{figure}[ht]
  \centering
  \includegraphics[width=\textwidth]{fig_res_sensitivity.pdf}
  \caption{Global sensitivity around \textsf{C1} (primary evaluator).
  \textbf{(a)}~Morris $\mu^*$ of O17p with bootstrap intervals (15 factors, $r=20$ trajectories, 10 runs per point
  with common random numbers). Red ticks: $\mu^*$ of the split-half noise contrast.
  \textbf{(b)--(e)}~Sobol' total indices $S_T$ (bars, bootstrap intervals), first-order indices $S_1$ (open
  circles) and noise floor (red ticks) for O17p, persistent O13, $\rho_{\rm pk}$ and $f_0$. Saltelli design with
  $N=256$ and 11 factors. Dark bars meet the relevance criterion $S_T-S_T^{\rm noise}>0.05$ and
  $S_T-\mathrm{CI}>S_T^{\rm noise}$.}
  \label{fig:res-sensitivity}
\end{figure}

\subsection{Seeding density, lattice size and founders}
\label{app:res-p4p6}

\paragraph{Seeding density and lattice size (P4).}
The plane combines $L=15$, 20 and 30 with $N_0=20$--800 animals of age 4 weeks placed at random, plus four founding
pairs of age 7 weeks in one cell, with 30 runs per point and the other parameters of \textsf{C1}
(Fig.~\ref{fig:res-p4}). The peak density is monotone in $N_0/L^2$ (Spearman 0.998) and rises from 0.07 to 0.94;
points with the same density have nearly the same peak density whatever $L$, which supports, without testing, that
the lattice size adds nothing at fixed density (the fixed-density stage of Section~\ref{sec:res-robustness} tests
it directly). O17p reaches 0.8 for $0.22\le N_0/L^2\le0.89$. At low densities the persistent class and the failure
of the late cohorts are lost and some colonies survive; at the highest density the peak approaches the capacity and
O4 fails. With four founding pairs the peak density is 0.07 ($L=30$) to 0.17 ($L=15$), and with mating in the Moore
neighbourhood 0.18 at $L=30$; no founder run is Calhoun-like.

\begin{figure}[ht]
  \centering
  \includegraphics[width=\textwidth]{fig_res_p4.pdf}
  \caption{Seeding density and lattice size (P4; 30 runs per point). \textbf{(a)}~Peak density, \textbf{(b)}~empty
  fraction at the peak and \textbf{(c)}~O17p (95\% Wilson intervals) against the seeding density $N_0/L^2$, for
  three lattice sizes with random seeding (filled markers) and four founding pairs in one cell (open markers); the
  triangles pointing down use mating in the Moore neighbourhood. The vertical dashed line is the default initial
  condition ($N_0=400$, $L=30$); horizontal dotted lines mark the values of Universe~25 ($\rho_{\rm pk}=0.57$,
  about 20\% of empty nest boxes) and the 0.8 threshold.}
  \label{fig:res-p4}
\end{figure}

\paragraph{Founder protocol (P6).}
Four pairs of 7-week-old animals start in one cell of the $30\times30$ lattice, with mating in the Moore
neighbourhood, $T=900$ and 30 runs per point. The plane varies the initial social state of the founders $s_0$ (0.2,
0.3, 0.5; a value below one represents adaptation to the new enclosure), the number of movement sub-steps per week
$k$ (1, 2) and the longevity $a_{\max}$ (120, 200; with $a_{\max}=200$ the fertile window is extended to 133 weeks).
Table~\ref{tab:res-p6} lists the outcomes. The latency pattern O1 passes in up to 0.67 of the runs, but O1--O3
jointly with O17p pass in at most 0.43 [0.27, 0.61] ($s_0=0.2$, $k=1$, $a_{\max}=200$). The best points need
$a_{\max}=200$, outside the preset; with $a_{\max}=120$ no point exceeds 0.13. The peak density stays at 0.16--0.27.

\begin{table}[ht]
  \centering
  \caption{Founder protocol (P6): four pairs in one cell with mating in the Moore neighbourhood, 30 runs per point.
  $s_0$ is the initial social state of the founders and $k$ the number of movement sub-steps per week. O1--O3 is
  the fraction of runs that pass the latency, growth and deceleration patterns jointly; O17p with 95\% Wilson
  interval.}
  \label{tab:res-p6}
  \small
  \begin{tabular}{@{}cccccccccc@{}}
    \toprule
    $s_0$ & $k$ & $a_{\max}$ & O1 & O1--O3 & O17p & O13p & $P_{\rm ext}$ & $\rho_{\rm pk}$ & $f_0$ \\
    \midrule
    0.2 & 1 & 120 & 0.53 & 0.47 & 0.13 [0.05, 0.30] & 0.64 & 0.63 & 0.18 & 0.43 \\
0.2 & 1 & 200 & 0.67 & 0.63 & 0.43 [0.27, 0.61] & 0.80 & 0.93 & 0.23 & 0.33 \\
0.2 & 2 & 120 & 0.57 & 0.23 & 0.13 [0.05, 0.30] & 0.54 & 1.00 & 0.16 & 0.51 \\
0.2 & 2 & 200 & 0.63 & 0.43 & 0.23 [0.12, 0.41] & 0.74 & 1.00 & 0.18 & 0.46 \\
0.3 & 1 & 120 & 0.47 & 0.47 & 0.10 [0.03, 0.26] & 0.72 & 0.67 & 0.20 & 0.38 \\
0.3 & 1 & 200 & 0.43 & 0.43 & 0.33 [0.19, 0.51] & 0.83 & 0.97 & 0.24 & 0.30 \\
0.3 & 2 & 120 & 0.57 & 0.40 & 0.17 [0.07, 0.34] & 0.54 & 1.00 & 0.21 & 0.37 \\
0.3 & 2 & 200 & 0.37 & 0.30 & 0.23 [0.12, 0.41] & 0.80 & 1.00 & 0.27 & 0.24 \\
0.5 & 1 & 120 & 0.00 & 0.00 & 0.00 [0.00, 0.11] & 0.60 & 0.53 & 0.20 & 0.38 \\
0.5 & 1 & 200 & 0.03 & 0.03 & 0.03 [0.01, 0.17] & 0.90 & 0.87 & 0.24 & 0.29 \\
0.5 & 2 & 120 & 0.10 & 0.10 & 0.03 [0.01, 0.17] & 0.80 & 1.00 & 0.26 & 0.24 \\
0.5 & 2 & 200 & 0.23 & 0.23 & 0.10 [0.03, 0.26] & 0.93 & 1.00 & 0.27 & 0.23 \\
\bottomrule

  \end{tabular}
\end{table}

\subsection{Transplant experiments}
\label{sec:res-transplant}

Each of the 60 replicates per arm runs a mother colony and samples it at several times (Fig.~\ref{fig:res-transplant},
Tables~\ref{tab:res-trans} and~\ref{tab:res-trans-s}): in phase B; at $1.75\,t'_{\rm pk}$ (phase D of the primary
evaluator); at $1.75\,t_{\rm pk}$ with $t_{\rm pk}=\arg\max N$ (the preregistered protocol); at week 95 in every arm,
which removes the confound between rule and sampling time; and, not preregistered, at the first week in which
four males and four females past the plasticity window have $0.2\le s\le0.5$ (the ``band''). Four males and four
females are moved to an empty lattice; success is reaching 80 individuals within 150 weeks. In the mixed phases the
animals of one sex are paired with healthy animals of the other. Arm $X$ shares its mother colonies and its
transplanted animals with \textsf{C1} and places them in an enclosure without R1 and R2 with $r=0.1$.
\begin{itemize}
  \item \emph{Healthy controls.} Animals taken in phase B found colonies in 77--98\% of the trials in every arm.
  \item \emph{Phase-D animals under R1.} Whenever R1 is present (\textsf{C1}, \textsf{C1}$'$, \textsf{C0},
    \textsf{C2} and $-$R2), phase-D animals never found a colony at any of the three sampling times, nor with
    healthy partners (0/120 at week 95 in \textsf{C1}, \textsf{C1}$'$, \textsf{C0} and $-$R2; 3/120 in \textsf{C2}).
    Their social state is 0 (median), with all eight below $10^{-3}$ in 55--93\% of the trials.
  \item \emph{Without R1.} At week 60 the same test gives 0/60 in the arms without R1 too, because their colonies
    are still at the peak and the animals are young but fully deteriorated; at $1.75\,t_{\rm pk}$ (weeks 126--131)
    10/60 recover alone in $-$R1 and 3/60 in $-$R1$-$R2, and with healthy partners 34/120 and 14/120 at week 95. In
    arm $X$ the rates are 51/60 at week 60, 12/60 at week 97 and 5/60 at week 95: the recovery of the same
    animals is a matter of the fertile time left.
  \item \emph{Intermediate animals.} Band animals (weeks 8--11, ages 12--15, $s=0.33$--0.40) found colonies in
    5/60 (\textsf{C1}), 2/60 (\textsf{C1}$'$), 0/60 (\textsf{C0}), 7/60 ($-$R2) and 8/60 (\textsf{C2}) under R1,
    against 50/60 without R1, 46/60 without R1 and R2 and 57/60 in arm $X$. With healthy partners the arms with R1
    reach 35--46/60 (band males with healthy females) and 7--32/60 (band females with healthy males); the asymmetry
    follows from the fecundity $s_f^2s_m$. The equal-time comparisons between arms (Fisher, BH at 10\%) separate the
    arms with R1 from those without it in the band phases and in the mixed phases at week 95, and nothing else.
  \item \emph{Descendants of the band.} In the colonies founded by band animals, the adult descendants have a mean
    social state of 0.28--0.45 in \textsf{C1}, 0.26--0.42 in \textsf{C1}$'$, 0.11--0.35 in \textsf{C0} and
    0.33--0.37 in arms $X$ and $-$R1: in a small uncrowded colony the offspring learn the competence of their
    parents and stay there, with or without R1. The lineage criterion of Section~\ref{sec:meanfield} is therefore
    confirmed only in its first part, the founding bottleneck.
\end{itemize}
This reproduces Marsden's observation, reported by Calhoun~\cite{Calhoun1973}: animals removed in mid-phase D failed
to reproduce ``even placing them with adequate sex partners'' raised without crowding. The model reproduces it only
with R1, and by construction. The evaluator rates O12 as PASS in every arm with R1 and in $-$R1$-$R2, and as
MARGINAL in $-$R1 and arm $X$ when sampled at $1.75\,t_{\rm pk}$.

\paragraph{Transplant grid.}
Over 12 points of the plane $(\alpha,\kappa)$, with $\alpha\in\{2,4,6,8\}$ and $\kappa\in\{0.1,0.15,0.2\}$ and 40
transplants per point and phase, no phase-D transplant founded a colony (0/480), nor did any mixed transplant
(0/960). Phase-B transplants succeed in 0.63--1.00 of the trials, less often with stronger attraction (mean 0.97 at
$\alpha=4$ and 0.71 at $\alpha=8$), and O12 passes at all 12 points.

\begin{figure}[ht]
  \centering
  \includegraphics[width=\textwidth]{fig_res_transplant.pdf}
  \caption{Transplant experiments (60 replicates per arm and phase). $P_T$ is the probability that four males and
  four females moved to an empty lattice reach 80 individuals within 150 weeks.
  \textbf{(a)}~Animals taken in phase B and in phase D at three sampling times: $1.75\,t'_{\rm pk}$ of each colony,
  $1.75\,t_{\rm pk}$ with $t_{\rm pk}=\arg\max N$, and week 95 in every arm.
  \textbf{(b)}~Adults past the plasticity window with $0.2\le s\le0.5$, alone and paired with healthy animals of the
  other sex.
  \textbf{(c)}~Mean social state of the eight transplanted animals per trial, at week 95 and in the band.
  Arm $X$ moves the animals of the \textsf{C1} arm, the same individuals, into an enclosure without R1 and R2 with
  individual recovery at $r=0.1$. Dotted lines in (a) mark the O12 thresholds (0.1 and 0.5).}
  \label{fig:res-transplant}
\end{figure}

\begin{table}[ht]
  \centering
  \caption{Transplants that found a colony, by arm and phase (successes/trials; 60 replicates per arm). Phase D is
  sampled at $1.75\,t'_{\rm pk}$ (median sampling week in parentheses), then at $1.75\,t_{\rm pk}$ with
  $t_{\rm pk}=\arg\max N$, and at week 95 in every arm; ``+healthy (95)'' pools the two mixed phases at week 95
  (h.: healthy partners of the other sex). The band samples eight adults past the
  plasticity window with $0.2\le s\le0.5$ at the first week in which they exist (median week in parentheses), alone
  or with healthy partners of the other sex.}
  \label{tab:res-trans}
  \scriptsize
  \setlength{\tabcolsep}{4pt}
  \begin{tabular}{@{}lcccccccc@{}}
    \toprule
    Arm & B & D & $1.75\,t_{\rm pk}$ & week 95 & +healthy (95) & band & band M + h.\ F & band F + h.\ M \\
    \midrule
    C1 & 52/60 & 0/60 (60) & 0/60 & 0/60 & 0/120 & 5/60 (9) & 37/60 & 17/60 \\
C1$'$ & 46/60 & 0/60 (63) & 0/60 & 0/60 & 0/120 & 2/60 (8) & 35/60 & 11/60 \\
C1 $-$ R1 & 58/60 & 0/60 (86) & 10/60 & 0/60 & 34/120 & 50/60 (9) & 57/60 & 56/60 \\
C1 $-$ R2 & 50/60 & 0/60 (65) & 0/60 & 0/60 & 0/120 & 7/60 (10) & 39/60 & 7/60 \\
C1 $-$ R1 $-$ R2 & 58/60 & 0/60 (82) & 3/60 & 1/60 & 14/120 & 46/60 (9) & 54/60 & 43/60 \\
$X$ (C1 animals) & 59/60 & 51/60 (60) & 12/60 & 5/60 & 41/120 & 57/60 (9) & 60/60 & 58/60 \\
C0 & 47/60 & 0/60 (58) & 0/60 & 0/60 & 0/120 & 0/60 (11) & 39/60 & 18/60 \\
C2 & 49/60 & 0/60 (110) & 0/60 & 0/60 & 3/120 & 8/60 (10) & 46/60 & 32/60 \\
\bottomrule

  \end{tabular}
\end{table}

\begin{table}[ht]
  \centering
  \caption{Social state and age of the transplanted animals, by arm and phase: median over trials of the per-trial
  median $s$ of the eight animals / maximum $s$ over all trials; median over trials of the youngest and oldest
  animal (weeks). Last column: mean social state of the adult descendants (age $\ge a_w$) in the colonies founded by
  band animals, pooled over the three band phases and weighted by their number (in parentheses).}
  \label{tab:res-trans-s}
  \scriptsize
  \setlength{\tabcolsep}{4pt}
  \begin{tabular}{@{}lcccccccc@{}}
    \toprule
    & \multicolumn{2}{c}{D ($1.75\,t'_{\rm pk}$)} & \multicolumn{2}{c}{week 95} & \multicolumn{2}{c}{band} & descendants \\
    \cmidrule(lr){2-3}\cmidrule(lr){4-5}\cmidrule(lr){6-7}
    Arm & $s$ med./max & ages & $s$ med./max & ages & $s$ med./max & ages & $\bar s$ ($n$) \\
    \midrule
    C1 & 0.000 / 0.36 & 37--64 & 0.000 / 0.30 & 65--80 & 0.358 / 0.50 & 13--13 & 0.31 (10261) \\
C1$'$ & 0.000 / 0.60 & 40--66 & 0.000 / 0.49 & 60--80 & 0.354 / 0.50 & 12--12 & 0.27 (8156) \\
C1 $-$ R1 & 0.003 / 0.37 & 46--78 & 0.004 / 0.46 & 56--80 & 0.361 / 0.50 & 13--13 & 0.36 (39369) \\
C1 $-$ R2 & 0.000 / 1.00 & 40--67 & 0.000 / 0.61 & 59--80 & 0.362 / 0.50 & 14--14 & 0.19 (10131) \\
C1 $-$ R1 $-$ R2 & 0.025 / 0.80 & 47--78 & 0.028 / 0.81 & 52--80 & 0.330 / 0.46 & 13--13 & 0.23 (31128) \\
$X$ (C1 animals) & 0.000 / 0.36 & 37--64 & 0.000 / 0.30 & 65--80 & 0.358 / 0.50 & 13--13 & 0.32 (38964) \\
C0 & 0.000 / 0.60 & 36--58 & 0.000 / 0.29 & 66--80 & 0.396 / 0.50 & 15--15 & 0.11 (7448) \\
C2 & 0.000 / 0.63 & 30--79 & 0.000 / 0.71 & 28--79 & 0.398 / 0.50 & 14--14 & 0.35 (17350) \\
\bottomrule

  \end{tabular}
\end{table}

\subsection{\textsf{C1} versus \textsf{C2}}
\label{app:res-c1c2}

Table~\ref{tab:res-QC} lists the confirmatory contrast; 13 of the 29 rows differ after Holm adjustment with the
primary evaluator, 12 with the preregistered evaluator (O4 and O13p differ only with the primary O4 and the fixed class
threshold). The 29 rows hold 27 distinct tests: O6, O7 and A2 (recovery of natality) are decided by the same four \textsf{C2}
colonies, which survive to $T$ and recover natality, and in the production runs O17 and O17p differ (0.725 against 0.705). Nine rows are constant in both
arms and enter the Holm sequence with $p=1$, so the correction over 29 is conservative. Besides the losses in O8 and O9, natality collapse weakens in \textsf{C2}: O5 passes in 0.815 against
0.98, because its births are spread over a longer growth phase ($t'_{\rm pk}=62$ weeks, $t_{B99}=86$). Four of its
200 colonies persist to $T$ (A2 in 0.02). The decline is as long as in \textsf{C1}, $(T_{\rm ext}-t_{\rm pk})/t_{\rm
pk}\approx1.9$--2.1, so O8 fails because the mean age no longer rises at the required rate. The isolated class is
slightly smaller in \textsf{C2} ($\Phi^{\rm pers}_{\rm BO}=0.13$ against 0.17), the peak higher (0.42 against 0.36),
the empty fraction smaller (0.11 against 0.16) and extinction later (Kaplan--Meier median 218 against 169 weeks).

Fig.~\ref{fig:res-cohorts} shows the cohort structure behind O13b. In \textsf{C2} the isolated class is born later
than the crowded class in every run (median birth week 22 against 11), whereas in \textsf{C1} the two classes are
born together (10 against 11 weeks; the isolated class is later in 4\% of the runs): the isolated animals of
\textsf{C1} are refugees of the early boom, of the same cohorts as the crowded ones. In both candidates births are
front-loaded: half of the births of \textsf{C1} occur before $0.36\,t'_{\rm pk}$ and 93\% before $0.8\,t'_{\rm pk}$,
the start of the window used by the strict O13b criterion, which therefore cannot pass.

\begin{figure}[ht]
  \centering
  \includegraphics[width=0.8\textwidth]{fig_res_cohorts.pdf}
  \caption{Cohort structure of the deteriorated classes. \textbf{(a)}~Per-run median birth week of the persistent
  isolated class against that of the crowded class, at the time $t^*$ used by O13b, for \textsf{C1} and
  \textsf{C2} ($n=200$ each, jittered); the dashed line is the diagonal. \textbf{(b)}~Cumulative fraction of births
  against time in units of $t'_{\rm pk}$ for the \textsf{C1} reference ensemble (median and 10--90\% band over
  runs); the dotted line marks the start of the O13b window.}
  \label{fig:res-cohorts}
\end{figure}

\begin{table}[ht]
  \centering
  \caption{Confirmatory contrast between \textsf{C1} and \textsf{C2} ($n=200$ each, $T=700$; primary evaluator).
  Binary patterns are compared with two-sided Fisher tests, continuous magnitudes with Mann--Whitney tests, and
  $T_{\rm ext}$ with a log-rank test with censoring. Holm correction is applied over the 29 comparisons. Entries
  are fractions of passing runs, or medians, with the number of applicable runs in parentheses. The last column
  gives the two fractions and the verdict with the preregistered evaluator. Omitted rows show no difference: O10, O13
  and O14 pass in every run of both arms; O15, A1, A3 and A4 are identical in both arms; O8b, O11b and strict O13b
  are constant.}
  \label{tab:res-QC}
  \footnotesize
  \setlength{\tabcolsep}{4pt}
  \begin{tabular}{@{}lccllcl@{}}
    \toprule
    Pattern / magnitude & \textsf{C1} & \textsf{C2} & $p$ & $p_{\rm Holm}$ & Differs & prereg.\ evaluator \\
    \midrule
    O4 & 1.000 (200) & 0.900 (200) & $1.2\times10^{-6}$ & $2.2\times10^{-5}$ & \textbf{yes} & 1.00 / 1.00 no \\
O5 & 0.980 (200) & 0.815 (200) & $2.3\times10^{-8}$ & $4.7\times10^{-7}$ & \textbf{yes} & 1.00 / 0.90 yes \\
O6 & 1.000 (200) & 0.980 (200) & 0.12 & 1.00 & no & 1.00 / 0.98 no \\
O7 & 1.000 (200) & 0.980 (200) & 0.12 & 1.00 & no & 1.00 / 0.98 no \\
O8 & 0.645 (200) & 0.000 (196) & $<10^{-16}$ & $<10^{-16}$ & \textbf{yes} & 0.65 / 0.00 yes \\
O9 & 0.400 (200) & 0.300 (200) & 0.0462 & 0.74 & no & 0.40 / 0.30 no \\
O11 & 1.000 (200) & 0.995 (197) & 0.50 & 1.00 & no & 0.99 / 0.99 no \\
O13p & 1.000 (200) & 0.945 (200) & $8.5\times10^{-4}$ & 0.014 & \textbf{yes} & 1.00 / 1.00 no \\
O15b & 0.650 (200) & 0.452 (197) & $8.2\times10^{-5}$ & 0.0015 & \textbf{yes} & 0.65 / 0.45 yes \\
O17 & 0.980 (200) & 0.725 (200) & $4.0\times10^{-14}$ & $8.7\times10^{-13}$ & \textbf{yes} & 0.99 / 0.88 yes \\
O17p & 0.980 (200) & 0.705 (200) & $1.6\times10^{-15}$ & $3.8\times10^{-14}$ & \textbf{yes} & 0.99 / 0.88 yes \\
O13b (PASS or MARG.) & 0.000 (200) & 1.000 (200) & $<10^{-16}$ & $<10^{-16}$ & \textbf{yes} & 0.00 / 1.00 yes \\
A2 & 0.000 (200) & 0.020 (200) & 0.12 & 1.00 & no & 0.00 / 0.02 no \\
$\rho_{\rm pk}$ (median) & 0.361 (200) & 0.42 (200) & $<10^{-16}$ & $<10^{-16}$ & \textbf{yes} & yes \\
$f_0$ (median) & 0.163 (200) & 0.113 (200) & $<10^{-16}$ & $<10^{-16}$ & \textbf{yes} & yes \\
$\Phi^{\rm pers}_{\rm BO}$ (median) & 0.174 (200) & 0.133 (200) & $<10^{-16}$ & $<10^{-16}$ & \textbf{yes} & yes \\
$(T_{\rm ext}-t_{\rm pk})/t_{\rm pk}$ (extinct) & 2.06 (200) & 1.8 (196) & $3.9\times10^{-14}$ & $8.7\times10^{-13}$ & \textbf{yes} & yes \\
$T_{\rm ext}-t_{\rm last}$ (extinct) & 110 (200) & 110 (196) & 0.37 & 1.00 & no & no \\
$T_{\rm ext}$ (KM median; log-rank) & 169 (200) & 218 (200) & $<10^{-16}$ & $<10^{-16}$ & \textbf{yes} & yes \\
\bottomrule

  \end{tabular}
\end{table}

\subsection{Ensembles}
\label{app:res-ensembles}

Table~\ref{tab:res-ensembles} lists every ensemble of the production program with its size, horizon and base seed, so
that each number of the text can be traced to its source. Seeds are derived as
\texttt{SeedSequence([base\_seed, point, replicate])}; the base seeds of the production runs (25301--25336) differ from those
of the runs of the preregistered design, so the two sets of runs are independent.

\begin{table}[ht]
  \centering
  \caption{Ensembles of the production program (\texttt{results\_v2/f2b\_<stage>}). Points: design points or arms;
  per point: runs or transplant replicates; $T$: horizon in weeks; seed: base seed of the stage.}
  \label{tab:res-ensembles}
  \scriptsize
  \setlength{\tabcolsep}{4pt}
  \begin{tabular}{@{}llrrrrrl@{}}
    \toprule
    Stage & Content & points & per point & runs & $T$ & seed & used for \\
    \midrule
    \texttt{base} & reference ensemble: C1, C1$'$, C0, $\alpha=0$ & 4 & 200 & 800 & 600 & 25301 & production \\
\texttt{c1c2} & C1 and C2 & 2 & 200 & 400 & 700 & 25318 & QC \\
\texttt{ref} & short-lived and long-lived controls & 2 & 100 & 200 & 1000 & 25302 & Q3, controls \\
\texttt{nulos} & four null models & 4 & 200 & 800 & 600 & 25330 & discrimination \\
\texttt{abl1} & single ablations and variants & 9 & 120 & 1080 & 600 & 25303 & Q2 \\
\texttt{abl2} & double to quadruple ablations, $\gamma=0$ & 12 & 40 & 480 & 600 & 25304 & ablations \\
\texttt{trans} & 8 transplant arms, 11 phases & 8 & 60 & 480 & 300 & 25306 & Q1, O12 \\
\texttt{amax} & $a_{\max}\in\{90,\dots,150\}$ & 5 & 60 & 300 & 600 & 25305 & Q4 \\
\texttt{rob} & demographic and update variants & 16 & 200 & 3200 & 700 & 25331 & robustness \\
\texttt{ci} & non-synchronous initial condition & 4 & 200 & 800 & 700 & 25333 & robustness \\
\texttt{size} & $L\in\{15,\dots,60\}$ at fixed density & 5 & 60 & 300 & 600 & 25334 & size \\
\texttt{q7aw} & $a_w$ varied at $\gamma=0$ & 54 & 20 & 1080 & 600 & 25335 & Q7 \\
\texttt{p1} & plane $(r,w)\times\gamma$ & 112 & 30 & 3360 & 600 & 25309 & Q7, P1 \\
\texttt{p1aw} & plane $(r,a_w)$ & 56 & 30 & 1680 & 600 & 25310 & Q6 \\
\texttt{p2} & plane $(\alpha,\kappa)$, $p$ & 72 & 30 & 2160 & 600 & 25308 & Q5, P2 \\
\texttt{p5} & refuge plane & 50 & 30 & 1500 & 600 & 25311 & P5 \\
\texttt{lhs} & Latin hypercube 1 & 200 & 10 & 2000 & 400 & 25312 & global \\
\texttt{lhs2} & Latin hypercube 2 & 200 & 10 & 2000 & 400 & 25313 & global \\
\texttt{lhsx} & boundary points of the hypercubes & 18 & 40 & 720 & 400 & 25336 & global \\
\texttt{morris} & Morris screening & 320 & 10 & 3200 & 400 & 25307 & global \\
\texttt{sobol} & Saltelli design, 11 factors & 3328 & 6 & 19968 & 400 & 25314 & global \\
\texttt{robmu} & $\mu_{\rm bg}$ curve & 12 & 60 & 720 & 700 & 25332 & robustness \\
\texttt{p4} & seeding density and lattice & 24 & 30 & 720 & 600 & 25315 & P4 \\
\texttt{trgrid} & transplants on $(\alpha,\kappa)$ & 12 & 40 & 480 & 300 & 25316 & O12 \\
\texttt{refine} & 40 boundary points of the planes & 40 & 100 & 4000 & 600 & 25317 & planes \\
\texttt{p6} & founder protocol & 12 & 30 & 360 & 900 & 25319 & O1--O3 \\
\bottomrule

  \end{tabular}
\end{table}

\subsection{Deviations from the preregistration}
\label{app:res-deviations}

\begin{table}[ht]
\centering
\caption{Deviations from the preregistered plan (\texttt{docs/preregistered\_plan.md}, dated 1 October 2026, 22:02 UTC,
before the production runs) and post hoc
analyses. ``Before production'' means after the calibration but before the production runs; ``post hoc'' means
after the runs had been seen. The last block lists the changes that are not preregistered but were fixed, with their decision rules, before
the production runs (\texttt{docs/production\_plan.md}, 2 October 2026, 05:34 UTC).}
\label{tab:deviations}
\scriptsize
\setlength{\tabcolsep}{3pt}
\begin{tabular}{@{}p{4.9cm}p{2.6cm}p{4.3cm}p{3.7cm}@{}}
\toprule
Change & When & Affects & Treatment \\
\midrule
O13p (persistent isolated class) defined and made part of the primary outcome, together with R7 & before
production, after seeing that \textsf{C0} has no persistent class & O17p, Q2, necessity of R7 & declared as
circular for R7 (Section~\ref{sec:discussion}) \\
O6: rebound count also detects slow regrowth still rising at censoring & post hoc (2 Oct., 02:44) & one
\textsf{C1} run of 200 under the preregistered design; none in the production runs & primary evaluator \\
O13b: reading of the window against the median birth week (Fig.~\ref{fig:res-cohorts}) & post hoc & the
interpretation of O13b in \textsf{C1} and \textsf{C2} & labelled post hoc \\
Robust spread of $T_{\rm ext}$ (Kaplan--Meier interquartile range over median) & post hoc (the plan
had $\mathrm{CV}\le0.10$ for Q3) & the claims about the spread of the extinction time & labelled post hoc; the
equivalence of \textsf{C1}$'$ with the long-lived control found under the preregistered design does not hold in the production runs \\
Ensembles of \textsf{C1}$'$ and the \textsf{C1}--\textsf{C1}$'$ contrast & post hoc & minimality, the role of
AV & labelled post hoc \\
Subsets of the hypercubes (nested conditions on $c_R$, $m_c$, $f_Rc_R$, $\pi$) & post hoc & interpretation of the
Calhoun-like fraction & descriptive only \\
Re-scoring with other class thresholds $\theta$ and ensemble thresholds $\tau$ (Appendix~\ref{app:res-threshold}) &
post hoc & conditionality of minimality & descriptive only \\
Cure model for the extinction time & post hoc & survival summaries & no test reported \\
Holm $p$-value of Q6 computed from the binomial likelihood & as planned & Q6 & the quasi-binomial value enters the
Holm sequence; no decision changes \\
\midrule
\multicolumn{4}{@{}l}{\emph{Not preregistered, fixed before the production runs}}\\
Hard refuge capacity enforced strictly (pups and rejected movers counted); the single-pass admission kept as a variant &
before the production runs & every stage with refuges; all stages run with seeds different from those of the preregistered design & both admission rules reported side by side
(no change in O17p) \\
Primary evaluator redefined: peak time at the start of the plateau ($t'_{\rm pk}$); end of natality in O5 at 99\%
of the births; O4 also requires free space at the peak ($f_0\ge0.10$, an informal target of the calibration);
crowded class $m>8$, independent of $m_c$; O10 on the adult social state; O11 on the second half of the births, not
applicable when the late cohorts are empty; NA counts as failure in O17p & fixed after exploratory runs with the primary evaluator and before the production runs & every pattern
normalised by $t_{\rm pk}$, O4, O5, O10, O11, O13, O13p, O17p in every stage; the verdicts of Q1 (prediction), Q2
(null) and Q6 (prediction) & the full preregistered evaluator, and each option changed separately, reported next to the
primary one (Appendix~\ref{app:res-rescoring}) \\
O5 birth reference $t_{B99}$ instead of the last sustained birth $t_{\rm last}$ & after exploratory runs had
shown that $t_{\rm last}$ with $t'_{\rm pk}$ fails O5 in most \textsf{C1} runs (0.48 in the production runs,
Table~\ref{tab:res-rescoring}); a choice informed by the result & Q2: it is the only option on which the rejection
of its null hypothesis depends (with the preregistered evaluator 5 of 7 null hypotheses are rejected and Q2 is not;
with the primary one, 6 of 7) & both birth references reported for every arm; the count with the preregistered
evaluator given next to the primary one (Section~\ref{sec:res-Q}) \\
Confirmatory test of the equivalence of the exponents of $r$ and $a_w$ at $\gamma=0$ (TOST, margin $[2/3,3/2]$,
stage $a_w$ varied) & before the production runs & Q7 with $a_w$ varied & the design cannot carry it out as inference
(51 of 54 points at $P_{\rm ext}\in\{0,1\}$, separated logistic); replaced by a non-parametric summary per window
(Table~\ref{tab:res-q7aw}); the planned estimates are kept, labelled, in \texttt{results\_numbers.json} \\
O8 (senescent decline) of \textsf{C1}: 0.71 under the preregistered design, 0.62 in the production runs (0.64 with the
single-pass admission and the seeds of the production runs) & observed after the production runs & a desirable pattern; no
decision & reported; not a conclusion \\
Transplant of Q1 taken at $1.75\,t'_{\rm pk}$ instead of $1.75\,t_{\rm pk}$; sampling at a common week added &
before the production runs & Q1, O12 & both times reported; comparisons between arms only at equal time \\
Sobol' design with 11 factors (the 8 selected by Morris plus $w$, $r$, $a_w$) & before the production runs & global
sensitivity & indices conditional on the 4 fixed factors \\
Additional stages: null models, demographic and update variants, non-synchronous initial condition, transplants of
intermediate animals, extra replicates at the hypercube boundary, $a_w$ varied at $\gamma=0$ (additional confirmatory test), lattice size & before the production runs & robustness, discrimination, Q1, Q7, feasible region & decision rules fixed before the runs \\
\bottomrule
\end{tabular}
\end{table}
